% 本文件由 LaTeX 排版 Agent 根据有效 translation.json 自动生成。
% author=Ephraim the Syrian; work_id=3702; title=尼西比斯赞美诗

\chapter{尼西比斯赞美诗}

% structure: p0001 source=h1 target=omitted reason=page-title-already-rendered-by-parent

1．\Place{尼西比斯}之围 (I-III)\newline{}2．\Place{波斯}入侵 (IV-XII)\newline{}3．\Place{尼西比斯}主教 (XIII-XVI)\newline{}4．\Person{亚伯拉罕}他们的继任者 (XVII-XXI)\newline{}5．论撒殚与死亡 (XXXV-XLII)\newline{}6．论撒殚与死亡 (LII-LXVIII)\par

\SourceRecord{Translated by J.T. Sarsfield Stopford.；Nicene and Post-Nicene Fathers, Second Series；Vol. 13.；Edited by Philip Schaff and Henry Wace.；Buffalo, NY: Christian Literature Publishing Co.,；1890.；<http://www.newadvent.org/fathers/3702.htm>.}

% structure: p0001 source=h1 target=section reason=page-title

\section{\WorkTitle{尼西比赞美诗}}

% structure: p0002 source=h2 target=subsection reason=relative-heading-rank; base=h2

\subsection{赞美诗第1首}

1. 仁慈的天主，祢曾安慰了诺厄，他也安慰了祢的仁慈。他献上祭品，止住了洪水；他呈上礼物，领受了恩许。他以祈祷和馨香平息了祢的怒火；祢以誓言和彩虹向他施恩；好使洪水若试图伤害大地，彩虹便伸展对阵，驱逐洪水，安慰大地。祢既已宣誓和平，求祢持守它，愿祢的彩虹与祢的怒气抗衡！\par

\Emph{启}：向洪水张开祢的弓，因为看，它已掀起波涛攻击我们的城墙！\par

2. 主啊！在启示中已宣告：诺厄所洒的那卑微的血，完全遏制了祢对万代的怒气；那么，祢独生子的血将何等更有力量，以致洒它就能遏制我们的洪水！因为看，那些卑微的祭品——诺厄所献、并\Emph{藉它们}遏制了祢怒气的——\Emph{不过是}祂的奥迹，才获得功效。求祢被我祭台上的礼物所平息，为我止住这致命的洪水。如此，祢的两个标记将带来救恩：于我，是祢的十字架；于诺厄，是祢的彩虹！祢的十字架将劈开水海；祢的彩虹将止住雨洪。\par

3. 看！所有的波涛都困扰我，而祢却赐予了方舟\Emph{更多的}恩宠：因为只有波浪环绕它，而土丘、兵器与波浪却环绕我。它曾是祢的宝库，我却成了债务的仓库：祢以爱制伏了它的波浪，我在祢的怒中却孤零零地留在兵器中间；洪水托起了它，河流却威胁我。啊，那方舟的舵手，求祢在旱地上作我的领航！祢使它安息在山巅的港湾，也求祢使我安息在我的城墙的港湾！\par

4. 正义者重重地惩戒了我，但祂甚至在波浪中也爱了我。因为诺厄战胜了情欲的波涛，那波涛曾在他那一代淹没了舍特的子孙。因他的肉体反抗了加音的女儿们，他的战车便驰骋于波浪之上。因妇女未玷污他，他结合了野兽，在方舟中成对地联姻。橄榄以其油滋润面容，以其叶欣悦他们的容颜；而那条河的饮水素来令人喜悦，看，主啊，它的洪水却使我悲伤。\par

5. 我的罪污已被祢的公义看见，祢纯洁的双眼憎恶我。祢藉不洁之手聚集了众水，好为我做成罪的净化；并非要在其中给我施洗并洁净我，而是要在其中以恐惧惩戒我。因为波浪将激发祈祷，以洗去我的罪愆。那充满忏悔的景象，于我已成了一种圣洗。主啊，那本该淹没我的大海，求祢在其中以祢的仁慈淹没我的罪愆。在红海中，祢淹没了躯体；在这\Emph{海}中，求祢代之以淹没我的罪恶！\par

6. 祢以仁慈预备了一只方舟，好将一切残余保存其中。为不使祢在怒气中荒废大地，祢的怜悯造了一只木制的大地。祢将它们彼此倒空，又彼此归还。\Emph{但}我的土地已三次被充满又再次倒空；如今波浪反叛我，要淹没我内逃脱的残余。在方舟中祢保全了残余；求祢在我内，主啊，是的在我内保全一个酵母。方舟在山巅孕育；求祢使我在我的土地上生出被囚禁的人！\par

7. 主啊！求祢使我堡垒中的囚犯喜乐，祢曾以橄榄叶使那些囚犯喜乐！祢藉鸽子把医治送给那些在每一波浪中溺毙的病者；它进入并驱散了他们所有的痛苦。因为它的喜乐吞噬了他们的忧伤，哀悼在它的安慰中消失。如同将领鼓舞逃亡者，同样，\Emph{鸽子}也在被遗弃者中散布勇气。他们的眼睛品尝了和平的景象，他们的口急忙张开赞美祢。如同波浪中的橄榄叶，求祢拯救我，好使祢因我使我堡垒中的囚犯喜乐！\par

8. 洪水冲击，拍打我们的城墙：愿支撑万有的力量扶持它们！它不像沙土建筑那样倒塌，因为我并未将我的教义建立在沙土上：磐石要作我的根基，因为我已在祢的磐石上建立了我的信德；我信赖的隐秘根基，将支撑我的城墙。因为耶里哥的城墙倒塌了，是因她将信赖建立在沙土上。梅瑟在海中立了一道墙，因他的理解将\Emph{它}建立在磐石上。诺厄的根基在磐石上；那木制的居所在海中承载了它。\par

9. 请将\Emph{在我内的}灵魂与\Emph{在方舟内的}生物相比；代替在其中哀伤的诺厄，看，祢的祭台在哀伤痛悔。代替在其中已婚的妻子，看，我的贞女尚未婚嫁。代替\Emph{从}方舟走出的含，看，正义的工人，他们曾养育并穿戴了宗徒们。在我痛苦中，我的主，我言语癫狂；若我的话惹动祢，求祢不要责备我！祢曾使抱怨的顺服者沉默：求祢怜悯我，如同怜悯从前被沉默的人！\par

10. 在祢的怒气之前，祢造了一座避难所，万民却反叛了它。诺厄在安息中得享安息，使他的居所按他的名字赐下安息。祢关闭了门以拯救义人，祢开启了洪水以毁灭不洁者。诺厄站在\Emph{外面的}可畏的波浪与\Emph{里面的}毁灭的口之间：波浪颠簸他，口使他惊慌。祢为他与里面的众物和好，祢降服了他外面的敌人：祢迅速转变了他的苦难，因为对于祢，主啊，艰难的事也是容易的。\par

11. 请听并衡量我与诺厄的比较，虽然我的痛苦比他轻，但求祢的仁慈使我们的救援相同；因为看！我的儿女像他一样，站在忿怒者和毁灭者之间。求主在境内赐予平安，在我面前谦抑境外的人，并赐我双倍的胜利！既然杀戮者已使他的愤怒三倍，愿那三日之主向我显示三倍的仁慈！愿邪恶者不能胜过祢的慈爱：看他已两次攻击我，而祢三次胜过他！愿我的胜利传扬于普世，好使它在世上为祢赢得赞美！哦，第三日复活的主，在我们\Emph{第三}次危难中，求祢不要将我们交与死亡！\par

% structure: p0015 source=h2 target=subsection reason=relative-heading-rank; base=h2

\subsection{赞美诗二。}

1. 今日我们张开嘴，以献上感恩。打开缺口的人，已打开我儿女的口。感谢仁慈者，祂拯救了我们的\Emph{城}中民众，那时并未想到索取我们应欠的债务。当他们起来，那掳掠我们的人，在我们的救援中，万有都尝到了祢的恩慈。\par

\Emph{叠句}。愿从一切有口者，荣耀归于祢的恩宠！\par

2. 祂无墙垣拯救了我们，并教导我们祂是我们的墙垣；祂无君王拯救了我们，并使我们知道祂是我们的君王；祂在每一个和所有事物中拯救了我们，并显示祂是万有；祂以恩宠拯救了我们，又启示祂白白地怜悯和赐予生命。祂从每一个自夸者那里夺去他的夸耀，归于祂自己的恩宠。\par

3. 众口的声音，不足以赞美祢；因为看！当我们的光冒烟，行将熄灭时（既然一切都易如祢），忽然它苏醒发光！谁曾见过这两件奇事：对他的希望被截断，希望却萌发增长；哀恸的时刻变为喜讯？\par

4. 这是一个庆节，万节依赖于它：因为若愤怒掳掠了我们，看！我们的节庆也终止了。既然我们的和平得胜凯旋，看！我们的节庆响彻云霄。这蒙福之日支撑一切：城邑依赖于它，人民依赖于城邑，和平依赖于人民，万有依赖于和平。\par

5. 从这些缺口，祢增加了胜利。赞美归于三一的天主，从三个缺口升起；因为祂降临并修复它们，以祂抑制忿怒的慈悲。祂击打\Emph{敌人}，而敌人不明白祂在教导我们。祂教导境内的人，因祂以公义制造缺口；祂教导境外的人，因祂以良善修复它们。\par

6. 说话并光荣吧，我这日得救的众子；老人与少年、青年与少女、孩童与无辜者，还有你，教会，城市的母亲！因为老人已从被掳中获救，青年从酷刑中，婴孩从被摔碎中，妇女从羞辱中，教会从嘲弄中。\par

7. 祂以严酷来到我们这里；我们惧怕了一时：祂以温和来到，我们欢庆了一刻。祂转身离开我们片刻，我们便无止境地流浪；如同猛兽，被甜言与恐惧驯服，但若人转而\Emph{离开它}，它便反叛、迷失，在和平中变得凶残。\par

8. 祂惩罚我们，我们却不畏惧；祂拯救我们，我们却不羞愧：祂束缚我们，我们的誓愿增多；祂放宽\Emph{我们}，我们的罪行增多。当祂约束时，有盟约；当祂赐予喘息时，有迷失。祂虽知我们，仍屈尊建立我们。傍晚我们高举祂，早晨我们弃绝祂。当急需离开我们时，信德也离开了我们。\par

9. 祂藉缺口磨难我们，以惩罚我们的罪行；祂筑起土垒，由此以谦抑我们的骄傲。祂为海洋开缺口，由此以洗去我们的污秽。祂关闭我们，好使我们聚集在祂的圣殿里。祂关闭我们，我们便熄灭；祂释放我们，我们便迷失。我们如同羊毛，浸入各种颜色。\par

10. 我们知道，当尼尼微蒙福的子孙悔改时，并非因土垒悔改，也不是藉着水，也不是因缺口，也不是因弓；他们不是因弓弦的响声畏惧而悔改。他们听从微弱的声音；使幼童禁食，使青年贞洁，使君王谦卑。\par

11. 祢击打我们，我们称祢为义，因非偶然；祢拯救我们，我们感恩，因非我们配得。祢怜悯我们，并非因祢错以为我们会悔改。祢明知，当祢怜悯我们时，我们迷失。祢知我们犯了罪；祢知我们是罪人：祢怜悯我们时，已熟知我们那已有的和现在的过犯。\par

12. 衡量我们的悔改，使它超过我们的罪行！但天平两端并不等重；因为我们的罪行沉重而繁多，我们的悔改轻微。祂曾下令我们因债务被出卖：祂的慈悲成为我们的辩护人；本金与利息，我们用悔改献上的小钱偿还了。\par

13. 为那微小的\Emph{偿付}，祂豁免了我们一万塔冷通的债务。祂本应索取，以平息祂的公义；祂又被迫宽恕，以让祂的恩宠喜乐。我们给了祂眼眨间的泪水；祂在索取并收取微小时满足了公义；当为微小而宽恕众多时，祂使恩宠喜乐。\par

14. \Emph{却是}万般罪过，祂已赦免；万般口舌，不足以在祂的良善面前赞颂。祂赦免了我们，我们却不赦免；我们以相反的方式回报祂；将所犯的罪重新写下。\Quote{主，求祢赦免，}我们呼喊；\Quote{主，求祢报应，}我们祈祷：\Quote{赦免}确实，当我们做了错事；\Quote{报应}确实，当错事施加于我们。\par

15. 是的，不像那些外邦人，我们为生命劳苦。他们筑起土垒，我们却连声音也没有；他们打破围墙，我们却连锁链——我们内心脆弱的锁链也没有折断。\Emph{天主}弃绝了勤勉者，为了懈怠者的缘故；祂弃绝了外表的劳作，虽然祂被从内心弃绝。\par

16. 祂释放了能言者，却击打了哑默者；城墙被攻破，百姓受管教：祂宽待那能受苦的，击打了那不知受苦的。祂没有击打能感觉的灵魂，却击打了无感觉的石头，为要管束我们。出于祂的爱，祂宽待了我们的身体，却急忙击打我们的城墙。\par

17. 有谁曾见过，一个缺口成了一面镜子？两方朝它观看，它同时为外人和内人服务。他们如同用眼睛在其中看见了那拆毁又建造的大能；他们看见了那制造缺口又予以修补的。外人在其中看见了祂的威能，便离去，未等到傍晚；内人在其中看见了祂的援助，虽感恩却仍不满足。\par

18. 愿你们得救的日子，唤醒你们脱离怠惰！当城墙被攻破，当大象冲入，当标枪如雨，当人奋勇之时，那时有了一场为天上者所见的景象。不义在那里争战，慈悲在那里得胜；仁爱在地上得胜，守望者在天上欢呼。\par

19. 你们的敌人劳累了，\Emph{竭力}以诡计击打那环绕你们的城墙，即你们居民的保障。他劳累了，却无济于事；为使他不存希望，以为一旦攻破便能进入并掳掠我们，他便攻破了它，且不只一次；他蒙受羞辱，却仍不满足，直至三次，好叫他在三次中三次蒙羞。\par

20. 愿我的幸福，因天主的恩宠，也在你们中间倍增！正如在你们中间我的罪行众多，愿在你们中间我的果实也众多！正如在你们中间我年轻时犯了罪，愿在你们中间我的老年得到怜悯！藉你们儿女的口为你们的儿子祈祷，因我犯罪超出了我的能力，悔改却低于我的能力；我撒种多，收获少。\par

% structure: p0037 source=h2 target=subsection reason=relative-heading-rank; base=h2

\subsection{赞美诗 三。}

1. 坚固我们的听觉，免得它松懈游荡！因为若有人追问祂是谁、祂像什么，它便游荡。我们如何能描绘那与心智相似的存有之像在我们心中呢？其中没有任何有限之物，祂全部看见、全部听见；祂整个彷佛在说话；祂整个充满各种感官。\par

R. 愿赞颂归于那独一的存有，我们无法测度祂！\par

2. 祂的容貌无法辨认，以至于能被我们的理解所描绘：祂无需耳朵而听；祂无需口舌而说；祂无需双手而工作，祂无需眼睛而看。因我们的灵魂在如此者面前并未停止或止息；祂以仁慈披上了人性的样式，将我们聚集到祂的肖像中。\par

3. 让我们学习那存有如何是属神的，却显现为有形的；又如何是宁静的，却显现为震怒的。这些事是为我们的益处；那存有按我们的样式与我们相似，好使我们能变得与祂相似。有一位与祂相似，即那从祂发出的圣子，祂身上印有祂的肖像。\par

4. 啊，尼西比，聆听这些事，因为这些事是为你的缘故而写、而说的。对于你自己和他人，你在世上成了争竞和辩论的缘由。口舌对着你，你这被围困的，甚至口舌为你歌唱；当你得胜、得释放时，在你里面口舌打开了，为哀悼和感恩。\par

5. 你居民的祈祷，足以使你获救；并非因他们是义人，而是因他们是悔改的：他们怎样蒙羞，也就怎样急忙服从懲戒。在过犯和得胜中，他们有同样的份。谁的罪行大，愿他的果实也大；谁在披麻蒙灰中得胜，也在冠冕中得胜。\par

6. 你得救的日子，是万日之王。安息日倾覆了你的城墙，倾覆了忘恩负义者；圣子复活的日子，又重建了你的废墟；复活日按它的名高举了你，荣耀了它的称号。安息日放松了它的守望；因造成缺口，它归咎于自己。\par

7. 在撒玛利亚，饥荒盛行；但在你们中间，丰盛行。在撒玛利亚，忽然有丰盛闯入降临她；但在你们中间，忽然有海洋咆哮涌入你们。在她那里，一个孩子被吃，救了她性命；在你们这里，那永生且赋予生命的身体被吃；忽然祂拯救了他们，那被吃的\Emph{拯救了}吃者。\par

8. 我们知道那位当受赞颂者不愿意历代以来的苦难；虽然祂行了这些事，却是我们的过犯成为我们患难的缘由。无人能抱怨我们的造物主；反倒是祂有理由抱怨我们，我们犯了罪，迫使祂虽不愿发怒却发怒，虽不愿击打却击打。\par

9. 大地、葡萄树和橄榄树，都需要修剪。橄榄被压榨时，它的果实香气四溢；葡萄树被修剪时，它的葡萄美好；土壤被耕种时，它的收成丰盛。水被约束在渠道里，沙漠之地得以畅饮；铜、银和金，打磨后发光。\par

10. 既然人藉\Emph{惩诫}能使万物美好；而那轻视、拒绝惩诫的，为人憎恶，众人都反叛他；\Emph{那么}就让他藉所惩诫的，学习那惩诫他的；因为凡惩诫者\Emph{这样做}是为使人\Emph{因此}获益。凡惩诫自己仆人的，\Emph{这样做}是为拥有他们；良善的天主惩诫\Emph{祂的}仆人，是为使他们拥有自己。\par

11. 愿你们的痛苦成为训诫你们的书籍，因为那三次被围困的经历，足以成为你们随时默想其历史的书籍。因你们轻视了那两部使你们能读得生命的约书，祂便为你们写下了三部硬书，使你们在其中读到自己的惩罚。\par

12. 让我们藉那已发生的，避免那将要发生的；让我们藉那已来的，学会逃脱那将来的；让我们记住那已过的，以避开那未来的。因我们已忘记了第一次击打，第二次便落到我们身上；因我们忘记了第二次，第三次重重压在我们身上。谁还会再次忘记呢！\par

\SourceRecord{Translated by J.T. Sarsfield Stopford.；Nicene and Post-Nicene Fathers, Second Series；Vol. 13.；Edited by Philip Schaff and Henry Wace.；Buffalo, NY: Christian Literature Publishing Co.,；1890.；<http://www.newadvent.org/fathers/3702a.htm>.}

% structure: p0001 source=h1 target=section reason=page-title

\section{\WorkTitle{尼西比赞美诗}}

% structure: p0002 source=h2 target=subsection reason=relative-heading-rank; base=h2

\subsection{第四首赞美诗}

1. 我的天主，我要不停践踏你殿宇的门槛；我既拒绝了所有恩宠，就大胆地求，使我能满怀信心地领受。\Emph{应答：我们的希望，请作我们的壁垒！}\par

2. 因为，主啊，大地使一粒麦子倍增百倍，那么我的祈祷岂不更因你的恩宠而丰盈！\par

3. 因我儿女的呼声、他们的叹息和哀号，求你向我开启你的仁慈之门！因他们的声音，使那穿麻衣者的悲伤转为喜乐！\par

4. 哦首生者，你曾是个断奶的\Emph{孩子}，与那些\Place{纳匝肋}的被咒诅的儿子们相熟，求你垂听我的羔羊，它们看见了豺狼，因为看啊！它们在哀哭。\par

5. 因为，我主，羊群在田野里，若看见豺狼，就逃向牧人，躲避在他的牧杖之下，他便驱散那要吞噬它们的。\par

6. 你的羊群看见了豺狼，看啊！它在大声哀号。看它何等恐惧！愿你的十字架成为牧杖，赶走那要吞灭它的！\par

7. 请收纳我那些全然纯洁的小孩子们的呼声。是那日子之婴孩，能够平息，主啊，那\OriginalTerm{万古常存者}{Ancient of days}。\par

8. 那天，当婴孩降生于马槽中时，守望者降临宣告了和平——愿那\Emph{和平}临于我所有的街巷，赐予我所有的后裔。\par

9. 那百姓的七十二位长老，不足以弥补其破口。是那婴孩，玛利亚之子，从四方赐下了和平。\par

10. 主啊，求你怜悯我的儿女！在我的儿女身上，求你记念你的童年，你曾是个孩子！愿那些像你童年的人，因你的恩宠而得救！\par

11. 在羊群中，混杂着无辜者的哭声，以及羊的呼声，它们呼求众羊的牧人，从一切中拯救它们。\par

* * * * * * * * * * *\par

13. 有一种喜乐是痛苦，其中隐藏着苦难；有一种苦难是利益，它是新世界里喜乐的泉源。\par

14. 我的迫害者所得的快乐，其中隐藏着灾祸；因此我喜乐。我因他所受的苦，其中为我隐藏着幸福。\par

15. 谁不赞美那生了我们、并能再次从恶謡中生出喜讯之声的那一位呢！\par

16. 万有的医治者，你在我的病中眷顾了我！我无法给你你的医药的代价，因为它们无价。\par

17. 你的仁慈比你的药物更加丰盛：它们不能购买，而是白白赐予，是以眼泪换取的。\par

18. 主啊，一座孤城，君王远在，敌人逼近，若无仁慈的\Emph{援助}，怎能屹立？\par

19. 你时刻是港湾和避难所。当海水淹没我时，你的仁慈降下，拉我出来。愿你的帮助再次抓住我！\par

20. 请将你救恩的药物和援助的激情，施于我的苦难！你的标记能成为治愈万有的良药。\par

21. 我深受压迫，急忙投诉那折磨我的。我的主，愿你的仁慈从我的罪所调的杯中除去苦味。\par

22. 我四顾垂泪，因我孤苦。尽管我的首领和拯救者甚多，但拯救我的只有一位。\par

23. 主啊，我的少年人已逃跑出去，如同雏鸡被鹰追逐；看啊！他们藏在隐秘处；愿你的和平带他们归来！\par

24. 我摘葡萄者的声音，看啊！我耳朵想念它，因他们的声音已绝。愿它因喜讯而回荡，哦，你救恩中的蒙福者！\par

25. 我听见恐惧的声音在我的城楼上；当我的守卫在城墙边呐喊时。求你以和平之声止息它！\par

26. 我农夫们的喧嚷，将在墙外宣讲和平；我居民们的呼喊，将在墙内宣讲和平，使我内外都享有和平。\par

27. 主啊，求你终结这纯洁祭台的哀悼，以及你那贞洁的司铎的哀悼，他穿着哀衣，满披麻布！\par

28. 教会及其仆役要为你的救恩赞美；城池及其居民亦然。主啊，愿和平之声成为他们声音的赏报！\par

% structure: p0031 source=h2 target=subsection reason=relative-heading-rank; base=h2

\subsection{第五首赞美诗}

1. 求你以恩宠使人听闻你救恩的喜讯：因为耳朵已被造成一条通道；我们的心思被恐怖的消息践踏。\Emph{应答：赞颂归于你的胜利！荣耀归于你的主权！}\par

2. 以微薄的收益安慰那些因劳苦而受损的人：在获利之时，他们所得\Emph{虽然}\Emph{却是}亏损。\par

3. 显然他已站立，将义怒分施于大地：在愤怒中他分赐损失和利益。有人被他骤然推倒，有人被他骤然高举。\par

4. 为教导\Emph{我们}他能够以各种方式惩戒；当他看见迫害者在我眼前可怕时，他将我展现在我儿女面前，而我所爱的儿女却鞭打我。\par

5. 看啊！他教导我敬畏他自己，而非人：因为当无人击打我们时，他的愤怒突然下令，每个人便伸手鞭打自己。\par

6. 同样，那巴比伦人，当他自信并以为无人能击打他时，他击倒了所有君王；天主却使他亲手击倒自己。\par

7. 他的威严和心智突然一同疯狂：他撕裂并抛弃衣服；出去在荒野中游荡；他先驱逐自己，然后他的仆人驱逐他。\par

8. 他向所有被他掳掠并降伏的君王显示，他并非凭自己的能力得胜：那击倒他的能力，正是惩罚他们的能力。\par

9. 我主，我站立承受了我的拯救者的打击。你能以你的恩宠使我因击打者得益；你能以你的公义借\Emph{我的}帮助者惩罚我。\par

10. 那日，敌军胆敢上来攻打\Place{撒玛黎雅}；他们的丰裕与欢乐、财宝与产业，都抛下弃绝逃跑。他借其迫害者加冕了她。\par

11. 我所爱的加冕了我，我的拯救者医治了我。因我居民的罪债，我的帮助者责打了我；求你从你的葡萄园中给我喝那安慰之杯！\par

12. 五谷与葡萄树，我主，求祢借恩宠保全！愿农夫因摘葡萄者的葡萄而欢欣，愿葡萄园丁因农夫的谷物而喜乐！\par

13. 五谷与葡萄彼此相连。在田地里，收割者能使酒欢畅；在葡萄园中，修剪者亦得面包滋养。\par

14. 这两样事物有能力安慰我的忧愁；天主圣三更有能力超然安慰；我愿赞美祂，因我骤然借恩宠获得了救恩。\par

15. 但若一人生命借恩宠得以保全，却因丧失财物而离去抱怨，他便对怜悯他的那位的恩宠忘恩负义。\par

16. 祂按己意毁灭一物以代替另一物。祂毁灭财产而保全拥有者；祂毁灭我们的植物代替我们的性命。\par

17. 让我们畏惧抱怨，免得祂的义怒被激起；祂保全财产，却击打拥有者；好叫我们最终明白祂起初的仁慈。\par

18. 让我们学习应当抱怨谁。学习抱怨，不是抱怨惩治者，而是抱怨你自己的意志，它使你犯罪并受惩罚。\par

19. 让我们除去抱怨，转向祈祷；因为若拥有者死去，他的财产也随他而止；但他若存活，便寻求挽回损失。\par

20. 愿安慰倍增，以怜悯我的居民；愿剩余和残存之物在忿怒中安慰 \Emph{我们}；并使我们借残存之物遗忘我们荒废的哀伤！\par

21. 我主，求祢医治并增加祢忿怒所遗留的果实！它们在我眼中如同在瘟疫中幸存的病者。求祢使我借这些弱者遗忘众多痛苦！\par

22. 我主，当我说话时，我忆起这正是那月，花朵凋零并因病枯萎落下；愿它恢复康健，成为安慰！\par

23. 因为这些逃脱了带走他们弟兄的瘟疫。葡萄树虽无声，当它们所爱的众多树木被砍伐倒下时，便在它们面前哭泣。\par

24. 众植物啊，大地正在思念它们！根苗为农夫哭泣，也使他们哭泣。它们的美丽曾伸展并遮荫，却在顷刻间被撕裂。\par

25. 斧头临近并击打；也击打了农夫；打击 \Emph{是}落在树上，却使农夫受苦；每一柄击打的斧头，他都承受其疼痛。\par

% structure: p0057 source=h2 target=subsection reason=relative-heading-rank; base=h2

\subsection{第六颂歌}

1. 我要以我的情感奔向那位白白医治者。祂医治了我的忧伤，第一次和第二次；\Emph{祂}治好了第三次，祂将医治第四次。\Emph{应答：首生之子，求祢治愈我！}\par

2. 我主，我的儿子们饮了忿怒所调和的讯息，便醉了；他们冲向我的装饰，掠夺并丢弃我的饰物；他们撕裂且毫不顾惜我的衣裳和冠冕。\par

3. 他们揭开我，使我赤露。因我蒙羞少许，经过那第一次和第二次的剥夺；因我第三次蒙羞，看，他们已第四次剥去我的衣服。\par

4. 因为他们已夺取并带走我的衣裳、我的饰物和我的花园。我主，求祢垂顾那束在我祭台上的苦衣，怜悯我！我主，愿这苦衣成为我救恩的护胸甲！\par

5. 看，我主，祢并非借洁身者之手责罚我！因他的羞耻在他面前，他的凌辱在他身后；至于他的婚姻，姦淫都比它好。\par

6. 看，他的女儿是他的妻子，他的姊妹是他的配偶；他转而娶那生他的母亲 \Emph{为妻}！诸天惊异他如此激怒祢，却仍亨通。\par

7. 我主，虽我的罪恶众多，但我的过犯如此沉重，以至于祢将贞洁之女——贞洁之女们的母亲——交予污秽的亚述——污秽之女们的母亲吗？\par

8. 求祢拦阻他，使他不得前来向我摇头、用脚踩我、并因他的名声震动世界而欢喜；也勿使他再稍得高举！\par

9. 我主，我的儿子们看见了我们的赤露，是的，他们揭开我并哭泣。求祢在我的儿女们面前揭露我，他们因我的痛苦而痛苦，不要让那些毫无怜悯的可咒者嘲笑我！\par

10. 我的土地曾结出果实和佳美之物；葡萄园中的美物，田野里的丰盈。但我安然休息时，忿怒忽然临到我。\par

11. 农夫被掠夺，抢掠者堆积 \Emph{谷物}；他们所借来撒种的，这些人都毁灭了。因债务，他的饥饿或许仍不得满足，因他的面包被夺去。\par

12. 我主，农夫被掠夺，因他贷与大地；大地接受了存款，却交给了外人；她向农夫借了，却付给了抢掠者。\par

13. 对祢，我主，我属于祢，求祢妒爱我！那将我许配给祢的宗徒告诉我祢是妒爱的。因为妒爱如同墙壁，保护贞洁的妻子。\par

14. 三松曾掀起海涛，因他极其妒爱他的妻子，虽她大大玷污并与他对立。求祢保全祢的教会，因她除祢以外别无他人！\par

15. 凡不妒爱自己配偶的，便是轻视她。妒爱能显明内在的爱。祢被称为嫉邪的，为要向我显示祢的爱。\par

16. 女子的本性即是如此：软弱而轻率；是妒爱使她时刻在畏惧之下。祢被列于妒爱者之中，为要显明祢的关怀。\par

17. 每个人都曾是某物的主人，那物 \Emph{并不}属于他；每个人都曾出去收集他没有散播的东西。混乱之日，我已因我的罪行而为自己预备。\par

18. 他们怎能忍受这痛苦，劳作者和耕种者？在葡萄园丁面前，他们砍倒葡萄树，赶走农夫的羊群；他们收割并夺走了他的播种。\par

19. 他们曾套牛、撒种、耙地；他们曾犁地、栽种、培育。他们远远站着哭泣，空手而去。劳动归于劳作者，收获归于抢掠者。\par

20. 统治者，我主，在\Emph{祢的}忿怒中没有维持秩序。倘若他们愿意\Emph{它}，他们本可保持秩序，但我们的邪恶不容许。虽然忿怒已大大减退，忿怒却强迫\Emph{他们}掠夺。\par

21. 我往哪边看，寻找安慰，因我的种植园被夷平，我的财产被荒废？愿平安声音的讯息驱散我的忧愁！\par

22. 不要将我交出；免得人以为祢给了我休书，将我送去并赶出！愿他们不称我为，我主，被遗弃的和受羞辱的！\par

23. 我没有任何事，可在祢眼前记起，因我完全被轻视。请为我记起，我的天主，只有这一件：除祢以外，我没有将任何别的置于我面前！\par

24. 谁不为我哭泣，以声音和哀号？因为在\Emph{满月}之日以前，我是贞洁且戴冠的；而在满月之日以后，我被揭露并赤身露体。\par

25. 我内室中贞洁的女儿们，在田野中流浪；因为那使所有人沉醉的忿怒，使我的尊贵妇女被轻视。愿祢那赐予众人平安的慈悲，使这些蒙爱者重获尊荣！\par

26. 我年长的\Emph{女儿们}和年幼的，看！她们在祢面前哭泣；少女以她们的声音，年老者以她们的眼泪；我的贞女以她们的禁食，我的贞洁者以她们的苦衣！\par

27. 我的眼目向所有街道举目，看！它们荒凉了。百人中剩下十人，万人中剩下千人。赐予和平，用我居民的喧嚣充满我的街道！\par

28. 把外面的人带回，使里面的人喜乐！祢的恩宠大有能力，祢将它延伸于内外。愿祢恩宠的翅膀，将我的小鸡聚集在一起！\par

29. 愿我义人的祈祷，拯救我的逃亡者！不信者掠夺了我，信者支持了我。在信徒中使不信者蒙羞！\par

30. 在同一天，两个庆节合而为一：祢升天节和祢的胜利者节；那编织祢荣冠的庆节，以及祢仆人接受冠冕的纪念。\par

31. 求祢垂怜，因为为我们，\Emph{这些}庆节在同一天加倍了；而为我们，代替它们，\Emph{甚至}两个庆节合而为一，却从\Emph{凶}信的声音中受苦，从荒凉中哀悼！\par

32. 赐予我的节期平安！因为我的两个庆节都已止息；而代替欢乐，我的残余在节期中，战栗和荒凉处处迎接我。\par

33. 将我远处的人带回家，使我近处的人喜乐；在我们的土地中间，要宣讲喜乐的好消息；为换取平安，我将从每张嘴里献上赞美！\par

% structure: p0091 source=h2 target=subsection reason=relative-heading-rank; base=h2

\subsection{第七首。}

1. 忿怒前来责备那些在和平中营谋、欺骗和掠夺的贪婪者。在灾难中贪婪者发了财：看！他们的东西他们已散尽，不是他们的东西他们已聚敛。\Emph{叠句：赐予和平，圣子，于我们的土地！}\par

2. 二十年来，我的烦恼，如同枝子，我的救主！它们在整个冬天被抑制，但当发芽的季节来到，我的烦恼就发芽：随着我们的果实，我们的心成熟。\par

3. 尼散月是萌芽的时候：在其中\Emph{凶}信萌芽了。当我们的喜悦向我们涌来时，那时我们的灾祸也向我们涌来。在簸扬小麦的时候，来到了城市的簸扬。\par

4. 因为\Place{巴比伦}的三个兄弟并未逃避人所点燃的火，因为他们坚定；他们逃避了情欲，因为他们成全。\par

5. 得胜者的火，能将黑色的小山羊变为白色；虚妄之人的火，能将羔羊变为有斑点的豹。\par

6. 我的哀号将多么大，在任何警报中呼喊！我的义愤多么大，在每个凶信中成熟！我的收成多么大，每张嘴都毁灭！\par

7. 因我儿子的罪行，祢惩罚了我，在他们为我的解救而奋斗时。那解救我的百姓，却将惩罚带到我身上。你们约束你们的罪，看！我的惩罚就被约束了！\par

8. 在凶信中他们受苦；在忿怒时他们被折磨；在和平时期他们却忧愁；因为当每个人自由呼吸，而\Emph{所有}的人都对恩宠忘恩时，他们替每个人献上\Emph{感谢}。\par

9. 他们的苦衣\Emph{是}因我而谦卑；他们的灰烬\Emph{是}洒在我的苦难中；他们的祈祷\Emph{是}为我的胜利；他们的禁食\Emph{是}为我的解救：看！债务在我苦修者身上，罪责在我贵族身上。\par

10. 在每个时代，智者的愚昧都是大的；经师和长老嫉妒并杀死了那教师，他曾教导万民\Person{梅瑟}的法律。\par

11. 今世的智慧\Emph{是}一种招致亏损的财产：人若有一点愚昧，他的罪责\Emph{是}很小的；但人若有一点明智，他的不义就越过了限度。\par

12. 他们用言语建造，却以行为推翻；因为教师众多而愚昧，但审判者的口却兼有这两样，审判者和控告者。\par

\EditorialNote{第八首圣歌缺失，第九首圣歌的开头部分亦缺失。}\par

% structure: p0105 source=h2 target=subsection reason=relative-heading-rank; base=h2

\subsection{第九首。}

…我的苦难如同\Person{约伯}的。祢的公义救了他；愿祢的恩宠怜悯我！\par

2. 在这两件事上有益处：既不让义人因恳求而疲倦，也不让悖逆者增加过犯。\par

3. 祢与儿子们一同劳作，为责罚并帮助\Emph{他们}；为使父老们不致因鞭声而悲伤，他们使我得享平安。\par

4. 看，我主，我外面的树林，它们被砍倒了！看，我主，我内在的胸怀，它们\Emph{太}软弱，使我不能怀抱我所爱的人！\par

5. 他们用刀砍断了我外面的翅膀；再次有火在我内在的胸怀中点燃了全燔祭的馨香。\par

6. 拜日者杀死了我平原上的儿子们；那向\Person{巴耳}献祭的，在城中牺牲了我的公牛，我的羊与我的婴孩。\par

7. 在我的田野中有哀号；在厅堂中有哭泣；在葡萄园中有惊恐；在街市中有混乱。谁能为我足够？\par

8. 那恶者行诡诈，用他的话扰乱我，在内心挑起麻烦，致使我里面的部分完全如同外面的部分。\par

9. 我主，我以何颜面，求祢派遣一队圣者，来护卫我这充满不洁的胸怀？\par

10. 祢以自己的新酵母惩治了受造界。求使那诱骗人、使人谦卑的旧酵母，变得像新酵母一样！\par

11. 借着祢权能的显明奋斗，我们得以得胜；免得谬误为那些为祢奋斗者加冕，并以谄媚依附于\Emph{他们}！\par

12. 若我们审视自己的时代，它犹如我们的诡诈；——因为在真理的岁月里，我们行占卜，并且秘密使用巫术。\par

13. 若我审视那时代，它激起并将隐秘之事带到光明中，好使我们那披着真理外衣的诡诈蒙羞。\par

14. 诚然是真理胜过一切；——海洋虽含苦涩，也不能搅扰它，——因其本性纯洁。\par

15. 祢以智慧造了它，它揭露了我们的贪欲。——为使愚昧者归于虚无，而不受鼓励——\Emph{真理}扣留了冠冕。\par

16. 在祢赐我胜利的倾颓城墙上——忘恩者以祭献和奠酒回报祢，这些公然激怒祢。\par

17. 倘若那时曾献上祭献，——那么即便错觉也可能以为——我因这些而得救。\par

18. 借着众多的拯救，祢斥责了两件事：——偶像的错觉和术士的教训；——因为我主，我因祢而得救！\par

% structure: p0124 source=h2 target=subsection reason=relative-heading-rank; base=h2

\subsection{颂 10.}

1. 我的儿女被杀，我的女儿们\Emph{那些离弃我的}——她们的城墙倾覆，她们的孩子四散——她们的圣地被践踏。——\Emph{回应：愿祢的惩罚受赞美！}\par

2. 捕鸟者从我堡垒中捉住了我的鸽子——它们离弃巢穴，逃往洞穴——却落入网中。\par

3. 如同蜡在火前融化——我众子的身体就在炎热和我堡垒的干旱前融化消散。\par

4. 取代以往涌流的甘乳之溪——为我的儿子和幼童，奶断绝于乳儿，水断绝于断奶孩童。\par

5. 乳儿从母亲怀中坠落喘息——因不能吮吸，母亲也不能哺乳——他们吐出气息而死。\par

6. 祢的恩宠怎能遏制其涌流之溪，既然无法遏制其丰沛的奔流？\par

7. 为何祢的恩宠封闭了它的怜悯，收回了它的溪流，使那呼喊要一人润舌的百姓得不到？\par

8. 在他们与同胞之间有一深渊——如同那呼喊却无人应答润舌的富人。\par

9. 如同投入烈火之中，悲惨者被投入——干渴中的炎热，火焰吹起——燃向他们。\par

10. 他们的躯干因炎热融化消散——曾干渴者反过来使大地饮下他们身体散发的气味。\par

11. 那以干渴杀死居民的堡垒——也转而饮下因干渴融化的尸身流质。\par

12. 谁曾见过干渴燃烧的百姓——四周环绕水墙，却无法润舌！\par

13. 我的亲爱者确实受审如同索多玛的审判——我的儿女遭击打，如同索多玛的折磨——尽管那\Emph{仅仅}是一日之事。\par

14. 火刑虽\Emph{仅}一时，我主啊——在持久的干渴中，是缓慢的死亡，精妙的惩罚。\par

15. 我主啊，在我忧愁与痛苦之后——这是祢安慰我的最佳慰藉——祢加增了我的苦难。\par

16. 我所希望的药，却是注定的愁苦；——我所期待的包扎，却是残酷的灾祸——它正为我施行。\par

17. 我本希望逃离风暴——港中的风暴对我比海中的更糟。\par

18. 我愚昧地以为能抛锚逃离海湾——我的罪恶却将我掷回其中。\par

19. 我主啊，请看我的四肢：刀剑密集刺入——在我臂上留下\Emph{其}印记；枪矛的疤痕——植入我的肋旁！\par

20. 我眼中流泪，耳中听闻\Emph{恶}讯——口中有哀歌，心中有忧伤！——我主啊，不要再加给我！\par

% structure: p0145 source=h2 target=subsection reason=relative-heading-rank; base=h2

\subsection{颂 11.}

1. 祢的惩罚\Emph{如同}我们幼年的母亲：——她的责备满含怜悯，因祢阻止了孩子行愚事——他们变得明智！——\Emph{回应：愿光荣归于祢的公义！}\par

2. 让我们探究祢的公义；谁又能充足——衡量其助益？因它常使放荡者变得贞洁。——\par

3. 我主啊，祢的手常常使病者痊愈——因它暗中医治他们的疾病——是他们生命的泉源。\par

4. 祢公义的手指极其温柔——以爱与慈悲触摸那要得痊愈者的伤口。\par

5. 对明智者而言，她的切割极其温和而仁慈：——她的锐利疗法，因着大爱——消融了腐败的\Emph{部分}。\par

6. 对明辨者而言，她的愤怒极其受欢迎——但她的疗法被那以肢体痛苦为乐的愚人憎恨。\par

7. 她极其急切地包扎自己所造成的伤口——当她击打时便怜悯，好从两者之中——孕育医治。\par

8. 她的愤怒极其受欢迎，她的怒气令人愉快——她的苦涩甘甜，使苦物变甜——好使之可口。\par

9. 祢的宽容对粗心者是疏忽之因——祢的杖对懒惰者是获益之因——好使他们成为\Emph{如同}勤劳者。\par

10. 我们受苦之因是祢的公义——我们疏忽之因是祢的恩慈——因我们的理智已变得愚拙。\par

11. 法郎因祢的恩慈而硬心；——当灾祸止息时，他的残暴增强——他背弃了自己的承诺。\par

12. 公义回报了他，因他大大撒谎反对她——甚至反对她的自由姐妹\Emph{恩宠}；她再次阻止他——使他不再挑衅。\par

13. 我主啊，请斥责我的向导，因它像埃及一样不忠——我的祈祷作证，我不像她——因我未曾离弃祢的门。\par

14. 我主啊，愿祢的十字架——它屹立在我敞开的破口中——再次修补隐藏的\Emph{破口}；因那些外在的破口，内在的已将我撕裂！\par

15. 海洋冲破，推倒了那曾使我得胜的瞭望塔。—— 不义竟敢建立一座使我蒙羞的殿宇：它的奠祭令我窒息。\par

16. 我在城墙上祈祷时，我的迫害者听见了：—— 太阳和它的崇拜者因他们的术士而羞耻，—— 因我藉祢的十字架得胜了。\par

17. \Emph{所有}受造物都呐喊，当他们看见那争战时—— 真理与谎言在我残破的城墙上搏斗，并戴上冠冕得胜。\par

18. 真理的力量责罚了谎言：—— 在责罚中它感到了\Emph{真理}，并通过它\Emph{自己}的罪恶，赢得了真理的胜利。\par

19. 我深感惊恐；因为自从我得救以来，—— 那些尊贵且强有力者，曾献身于我的祭台，如今在我内建造了邱坛。\par

20. 我的七种感官，我主啊，即使它们像泪泉，我的眼泪仍\Emph{太}少—— 不足以哀悼我们的毁灭。\par

21. 那\Emph{曾}披麻蒙灰的街道呐喊—— 被那与在旷野里金牛犊面前\Emph{曾有}的嬉戏相仿的嬉戏所搅扰。\par

22. 毒物寻求并披戴百合的美丽—— 虽然它们的花蕾可以隐藏，并暗中伪装它—— 它却在它们苦涩的花朵中绽放。\par

% structure: p0168 source=h2 target=subsection reason=relative-heading-rank; base=h2

\subsection{赞歌十二}

1. 我在苦难中，要呼求那制服万有的权能—— 它能制服怒中的俘囚者，—— 正如它曾克服\OriginalTerm{军团}{Legion}。\Emph{答：光荣归于祂的恩宠！}\par

2. 邪恶者偿还了我，我的弟兄们，他并未向我借的债务：—— 善良的\Emph{天主}也同样偿还了我，我并未借给祂的仁慈。—— 来，并惊异这两件事吧！\par

3. 善良的\Emph{天主}已分开并给予：我的恶行归于祂的恩宠—— 我的过犯归于祂的公义；祂的仁慈已抹去我的恶行—— 祂的审判已报应我的过犯。\par

4. 罪恶极其愤怒，并在惊惧中停留—— 当她看见恩宠如何约束自由，好让她能够战胜过犯。\par

5. 燃烧吧，我主，并倾下祢的爱，爆发并倾出祢的怒火！—— 祢的怒火为毁灭，祢的爱为拯救—— 将俘囚者手中的俘虏解救出来！\par

6. 那些日子，恶者曾决意将我抛出—— 如同用机弦抛入灭亡，在其中善良的天主已束好并保存—— 我的灵魂在生命囊中。\par

7. 那些言语之人，从不静默，不断赞美—— 他们曾在波涛中保守我，扶持我使我不致跌倒，愿他们代替我赞美，我主啊！\par

8. 因为谁曾足够面对那包围他的仁慈的恩宠—— 使我足以赞美那环绕我的仁慈呢？\par

\SourceRecord{Translated by J.T. Sarsfield Stopford.；Nicene and Post-Nicene Fathers, Second Series；Vol. 13.；Edited by Philip Schaff and Henry Wace.；Buffalo, NY: Christian Literature Publishing Co.,；1890.；<http://www.newadvent.org/fathers/3702b.htm>.}

% structure: p0001 source=h1 target=section reason=page-title

\section{尼西比圣咏集}

% structure: p0002 source=h2 target=subsection reason=relative-heading-rank; base=h2

\subsection{第十三首圣咏}

\Emph{关于\Person{玛尔雅各伯}及其同伴。}\par

1. 三位杰出的司铎，仿效两大光明——彼此传承交接了宝座、手和羊群。——对于我们这些因那两位而极其仰慕的人，这最后一位完全是安慰。\Emph{应答：愿光荣归于祢，是祢拣选了他们！}\par

2. 祂创造了两个大光，为自己拣选了\Emph{这些}三个光明——并将它们置于已经过去的三个黑暗围城时期。——当\Emph{那}一对光明熄灭时，另一个全然照耀。\par

3. 这三位司铎是宝藏，他们以信德持守着——天主圣三的钥匙；他们为我们开启了三个门——他们每人以自己的钥匙解锁并开启了自己的门。\par

4. 在第一位时，门被开启，为那降临于我们的惩罚——在第二位时，门被开启，为那降临于我们的君权——在最后一位时，门被开启，为那为我们而来的好消息。\par

5. 在第一位时，门被开启，为两军之间的战争；——在第二位时，门被开启，为来自两方的君王——在最后一位时，门被开启，为来自双方的使者。\par

6. 在第一位时，门被开启，为因恶行而起的战争；——在第二位时，门被开启，为因争斗而来的君王——在最后一位时，门被开启，为因仁慈而来的使者。\par

7. 看！在这三次传承中，如同一个奥秘和图像——义怒被比作太阳；它始于第一位之下——它在第二位之下增强；它在最后一位之下沉落并熄灭。\par

8. 太阳也在三个区域中显示出三个形象：——其升起\Emph{是}锐利明亮的；其中午强烈且压倒一切——而如同烧尽的火把，其落下\Emph{是}和缓愉悦的。\par

9. 其升起虽小却\Emph{是}明亮的，当它来唤醒睡眠者时——其中午炎热且压倒一切，当它来使果实成熟时——其落下柔和且愉悦，当它达到圆满时。\par

10. 这位因誓愿而生的女儿是谁，令所有女子羡慕——她的传承如此进行，她的队伍如此众多——她的等级如此上升，她的教师如此卓越。\par

11. 这些比喻只属于亚巴郎的女儿——还是也属于你，哦，因誓愿而生的女儿，你的装饰随你的美丽而变？——因为正如你的时机，你的援助\Emph{也是如此}，正如你的援助，其施行者\Emph{也是如此}。\par

12. 按照她需要的尺度，满足她需要的供应来临。——她的父亲们如同她的出生；她的教师们如同她的理解力——她的培育如同她的成长；她的衣裳如同她的身量。\par

13. 恩宠\Emph{为她}称量并赐予这一切，如同在天平上——她将之放在\Emph{她的}天平中，为能从中获益——她将它们纳入传承，为能成就完美。\par

14. 在第一位的日子里，平安充满且平安消失——在第二位的日子里，君王降临且君王退去——但在最后一位的日子里，军队进攻且军队撤退。\par

15. 借着第一位，秩序进入，与他同来亦与他同去——借着第二位，那使我们教会喜乐的冠冕靠近又远远退去——但借着最后一位，曙光降临于我们，那是未被感恩接受的恩宠。\par

16. 对抗那最初的义怒，第一位的劳苦相争——对抗那正午的炎热，第二位的荫蔽站立——对抗那忘恩负义的平安，最后一位倍增警告。\par

17. 因为这土地的第一位入侵者\Emph{正是}第一位杰出的司铎——这土地的第二位入侵者，正是第二位仁慈的司铎：——但最后一位的祈祷，秘密地修补了我们的破口。\par

18. 尼西比建于水上，隐秘与公开之水：——活水在她内；一条高贵的河流在她外。外面的河流欺骗了她；里面的泉源拯救了她。\par

19. 第一位司铎是她的葡萄园丁；他使她的枝条长到天上。——看！他死后被埋葬在她内，变成了她怀中的果实：——因此当修剪者来时，那在她内里的果实保全了她。\par

20. 她修剪的时候到了；它进入并带走了她的葡萄园丁，——以致无人为她祈祷。她以她的机巧急忙行动——祂将她的葡萄园丁放在她怀中，使她通过她的葡萄园丁得救。\par

21. 要像尼西比一样有智慧，哦，尼西比的女儿们，——因为她将身体放在她内，它就变成了她外的城墙。——将活的身体放在你们内，使它能成为你们生命的城墙！\par

% structure: p0025 source=h2 target=subsection reason=relative-heading-rank; base=h2

\subsection{第十四首圣咏}

1. 在三位牧者之下，——有众多的牧人——那城中的一位母亲——在各地有女儿。——既然义怒摧毁了她的住所——和平将建立她的教堂。\Emph{应答：愿受赞颂是那拣选了那三位的主！}\par

2. 第一位慈爱的劳动——在忧苦中包扎了土地：——第二位司铎的饼和酒——医治了破碎的城市：——最后一位甘甜的言语——在苦难中甘甜了我们的苦涩。\par

3. 第一位以他的劳动耕种土地——他拔出了其中的荆棘和蒺藜：——第二位在她周围围上篱笆——他用得救的人为她作了篱笆：——最后一位打开了他主人的粮仓——并将他主人的话语播种在她内。\par

4. 第一位司铎藉着禁食——封闭了众人之口：——第二位司铎为俘虏——打开了钱包之口：——但最后一位刺穿了耳朵——并将生命的装饰品系在其中。\par

5. 亚郎从耳朵上摘下——耳环并造了一只牛犊。——那只无生命的牛犊在暗中——刺穿并屠杀了营中的人：——那些曾铸造它角的人，——它用角撕开了他们。\par

6. 但我们的第三位司铎——刺穿了心灵的耳朵——并将他所造的耳环系在\Emph{那里}——用那钉在十字架上的钉子，——他的主人就曾被钉其上——并给予了他的同伴生命。\par

7. 火生出了一个死亡之子——死亡以所有身体为食：——死亡之子超越了死亡——他以人的灵魂为食。——牛犊离开了它的饲料——因为人的思想成了它的食物。\par

8. 第一棵树带来了死亡，——恩宠却为它生了一个儿子。——啊，十字架，那棵树的子孙，——你与你的父亲争战！——那树是死亡的泉源，——十字架是生命的泉源。\par

9. 死亡所生的儿子，——众口张开咒骂他。——他吞吃身体和灵魂，——加增了他父亲的羞辱。——但十字架除去了——它父亲、那第一棵树的耻辱。\par

10. 这两个儿子，正如——生育他们的\Emph{两个}母亲。——火生出的牛犊，——火在百姓中把它烧尽：——恩宠所生的十字架，——将美好的恩赐分给\Emph{所有}受造物。\par

11. 我的舌头啊，请静默不言，关于十字架的历史——它们正迫切地\Emph{有待述说}！——因为我的心思忽然怀了孕，——看哪！生产的剧痛击打它：——它在这末后怀了这些，——它们争着要成为首生者。\par

12. 婴孩在腹中相争，——年长的急忙先出：——年幼的想要长子的名分，——他的手抓住他的脚跟，——他在出生时未能得到的，——却因一碗红豆汤得到了。\BibleRef{创 25:29-34}\par

13. 同样，这些后来的历史，——看哪！它们轻视以前的\Emph{那些}，——好让\Emph{自己}出来夺取长子的名分。——让我们述说我们祖先的历史，——因为看哪！十字架的历史\Emph{是}\Emph{所有}受造物中的首生者。\par

14. 因为若那没有起始的，——是万有中的首生者，——它的历史也是首生者，——因为它们比\Emph{所有}受造物更古老。——求你那些历史，我的主啊，——让位，让我们述说你的仆人们的事！\par

15. 在教导的等级中居首者，——他的口才\Emph{正如}他的等级一般；——次位者，等级第二，——他的诠释达到了他等级的高度；——末位者，等级第三，——他的口才伟大\Emph{正如}他本人。\par

16. 首位者以他的纯朴言语，——将奶喂给他的婴孩；——次位者以他的简单话语，——将食物给他的儿女；——第三位者以他的完美言论，——将干粮给他的\Emph{那些已是}完全年岁的人。\par

17. 她，教诲的女儿，——也一级一级地上升，——与她的师长和父亲们同行。——与首位同在时，她是幼童；——与次位同在时，她是纯朴\Emph{的少女}；——在第三位那里，她达到了完全的年岁。\par

18. 首位\Emph{待她}如同孩童，——爱她，教她敬畏；——次位待她如同少女，责备她又使她喜乐；——第三位待她如同完全受教者，——成了她的愉悦慰藉。\par

19. 甚至至高者对待雅各的女儿，——在她的童年\Emph{赐}以抚慰和杖责，——在她叛逆和成年时，——赐给她刀剑和法律的分；——按照她的训练和教导，——祂以温柔和愉悦临到她。\par

20. 第一位生了羊群，——他的胸怀承载了她的婴孩期；——第二位面容欢愉，——以歌声鼓舞、使她童年喜乐；——第三位面容庄重，——看哪！他在她年少时守护\Emph{她的}贞洁。\par

21. 第一位司铎生了\Emph{她}，——给她的婴孩期喂奶；——第二位司铎诠释，——给她的童年食物；——第三位司铎养育\Emph{她}，——给她的完全年岁干粮。\par

22. 那位\Emph{是}首位富有的父亲，——为她的童年积蓄财宝；——次位为她的成熟，——增添了旅途的供应；——第三位，那美好的橄榄树，——在她器皿中加添了油。\par

23. 当她来到富足者面前时，——她将显明首位的财宝；——当她来到救主面前时，——她将显明次位所救的人；——当她出去迎接新郎时，——她将显明她灯里的油。\par

24. 在赏报劳苦困乏者的那位面前，——她将献上首位的辛劳；——在喜爱\Emph{乐意}施与者的那位面前，——她将显明次位的施舍；——在审判教义的那位面前，——她将献上末位的言论。\par

25. 而我，这个罪人，曾努力成为——这三位门徒——当他们看见第三天的那位——祂已关闭了祂内室的门时，——愿这三位为我向\Emph{祂}祈求，——求祂为我将门\Emph{留开}片刻！\par

26. 愿罪人挤进去进入——欢喜又战兢地观看！——愿三位师傅呼唤——那一位门徒，靠\Emph{他们的}恩宠！——愿他在桌子下面拾起——那充满生命的碎屑！\par

% structure: p0052 source=h2 target=subsection reason=relative-heading-rank; base=h2

\subsection{第十五首赞美诗。}

1. 若头不正，——肢体或许会抱怨：——因为因着一个乖僻的头，——肢体的运行被引入歧途，——他们惯常将过错归咎于头。\Emph{应答：愿祂受赞美，祂拣选了你们，我们子民的荣耀！}\par

2. 若如今我们向那全然美好的——倾注我们的仇恨，——那么若我们是可恨的，\Emph{更}当如何！——是的，即使天主是良善的，——心怀苦毒的人也抱怨祂。\par

3. 要像头一样，你们肢体啊！——在他的纯洁中得安息，——在他的安宁中得愉悦，——在他的圣德中得名声，——在他的智慧中得学识！\par

4. 在他的温和中得辨识，——在他的庄重中得贞洁，——在他的贫穷中得慷慨！——既然他全然美好，——让我们也与他全然美好！\par

5. 看哪，他的言语和行为——是何等量度和平衡！——你们要注意，甚至他的脚步——也持守平安的尺度！——他竭尽全力勒住整个自我。\par

6. 他作主管理他的青春，——将它捆在贞洁的轭下：——他的肢体未曾被\Emph{情欲}引诱，——因为它们被置于杖下：——他的意志被制服了。\par

7. 因为他早已为他的等级预备好，——正如他在生活上预备好自己，——他稳固地奠立根基。——他在年少时就成为头，——那时他们立他为百姓的讲道者。\par

8. 他在讲道者中超群，——他在学者中有学问，——他在智慧人中明达：——他在兄弟中贞洁，——在亲密朋友中庄重。\par

9. 他在两个居所中——从\Emph{早年}起就是独修者，——因为他在身体内是圣洁的，——在住所内是独居的，——公开和私下他都贞洁。\par

10. 但是我们，我的兄弟们，——虽然我们已偏离了那些尺度，——并失去了那滋味，——成了自己的教师，——背离了那召叫我们的完美。\par

11. 然而那真理的尺度——在其器皿中保存自身：——\Emph{真理}拣选它，因她见它拣选了她，——她在其中保存了她的芬芳和滋味，——从起初直到末了。\par

12. 那贞洁而庄重的元首——既不愤怒也不严苛——也不像我们一样\Emph{确实}越轨——确立并持守了自己的尺度——并在自己的思想上套上了缰绳。\par

13. 他以自身树立了榜样——好使我们像他持守自己时代的尺度一样——也应当认识自己的时代。——我们对自己的时代成了陌生人——因为在明辨的时代，我们失去了智慧。\par

14. 起初，猛烈的风力——以其威力催逼果实——然后，太阳的威力接踵而至：——但当其威力消逝后——它的终点聚集了它的甘甜。\par

15. 但我们——先来者管教了我们——后来的也责备了我们——最后者又为我们增添了甘甜：——然而当品尝我们的时刻到来时——我们却是何等的寡淡无味。\par

16. 我们本已成熟——本应使孩童戒除放荡——引导他们走向庄重：——\Emph{但}我们的年迈反倒需要——像青年人一样被责备。\par

17. 因此，他仁慈地容忍，并未使用强制——为要增加对我们年迈的尊敬：——即使它不认识自己的等级——愿那认识它时代者受赞美！\par

18. 若有人说，对于群众——应当用强制和棍棒治理——如同对窃贼用恐惧——对抢劫者用威胁——对愚昧者公开羞辱。\par

19. 然而，若以元首为首——肢体们迅速\Emph{行动}为次——他们便会带动那第三的——于是整个身体从末端——将跟随他们。\par

20. 那些\Emph{是}次等的，轻视了那些\Emph{是}首等的——那些\Emph{是}第三等的，轻视了那些\Emph{是}第二等的：——等级彼此被藐视。——当这些内部的人互相轻视时——他们同样被外部的人践踏。\par

% structure: p0073 source=h2 target=subsection reason=relative-heading-rank; base=h2

\subsection{赞美诗 16.}

1. 这里\Emph{是}一面该受责备的镜子——若其清澈变得昏暗——因其本体上有斑点；——因为其上的污秽——成为观看者眼前的遮蔽。\Emph{R.}愿那擦亮我们镜子的受赞美！\par

2. 因为美丽在其中并未得到装饰——瑕疵也未在其中得以显现——它完全损害了美丽的事物——因为它们的美丽未能获得——作为其益处的装饰。\par

3. 瑕疵不能藉它根除——同样，装饰也不能藉它增多。——留下的瑕疵如同亏损——没有装饰是缺陷：——亏损与缺陷相遇。\par

4. 若我们的镜子\Emph{是}黑暗——那对可憎者完全是喜乐；——因为他们的瑕疵\Emph{是}不被责备：——但若它被擦亮而闪耀——被装饰的乃是我们自由。\par

5. 亏损中有双重的损失——对可憎者与对美好者——因为美好者未获冠冕——同样，可憎者未得装饰：——镜子平分了亏损。\par

6. 镜子从不——强制它所照之人：——\Emph{同样}，恩宠——紧随律法的义德而来——并不拥有律法的强制力。\par

7. 义德对于童年——是强制的装饰者——因为当人类处于童年时——她以强制装饰了他——却未剥夺他的自由。\par

8. 义德使用温和劝诱——以及棍棒\Emph{来对待}童年——当她击打他时，她唤醒了他；她的棍棒克制了任性，她的温和劝诱柔软了心志。\par

* * * * * * * * * * *\par

9. [若有人偏离福音，]即我们今日所装饰的，我的弟兄们，——转向另一福音，他便是孩童：——在理解力伟大的时代——他变得没有理解力。\par

10. 因为在成年人的层次——他降到了童年——并且喜爱奴仆的律法——这律法在他自信时击打他——在他欢乐时掌掴他。\par

11. 任何强制性的装饰——都不是真实的，而是借来的。——这在天主眼中是伟大的——即一个人应藉自己装饰自己：——因此祂取走了强制。\par

12. 因为正如出于\Emph{祂}的明智——在\Emph{其}适当的时代祂使用了强制——\Emph{同样}，出于祂的明智——祂在某个时代取走了它——那时温柔和蔼取代了强制。\par

13. 因为正如适合青年——应在棍棒下急行——\Emph{同样}，\Emph{是}在棍棒下——让智慧服役，让强制作她的主——是非常可憎的。\par

14. 因此，请看同样——天主如何安排了我的先后次序——在我所拥有的牧者们身上——在祂赐予我的教师身上——以及在祂算给我的父亲们身上！\par

15. 因为按他们的时代称量——\Emph{是}他们品性的帮助；——即在需要恐惧的人身上，恐惧；在需要振奋的人身上，振奋；在需要温柔的人身上，温柔。\par

16. 祂以尺度使我前进：——对我的童年，祂指派了惊惧；同样，对我的青年，恐惧——对我的\Emph{年龄}智慧和明智——祂指派并赐予了温柔。\par

17. 在童年层次的任性别中——\Emph{我的}导师对我来说是一种恐惧：——他的棍棒约束我远离放荡——远离恶作剧的是他的惊惧——远离纵欲的是对他的恐惧。\par

18. 祂赐给我另一位父亲以应对我的青年：——我身上尚存的稚气——对应他身上的严厉；我身上的成熟——对应他身上的\Emph{如同}温柔。\par

19. 当我从童年和青年的层次升起时——那\Emph{是}第一的惊惧消逝了——那\Emph{是}第二的恐惧消逝了——祂赐给我一位仁慈的牧者。\par

20. 看！为我的成年，他的食粮——为我的智慧，他的诠释；——为我的平安，他的温柔——为我的安宁，他的仁慈——为我的贞洁，他的庄重！\par

21. 愿祂受赞美，祂如同在天平上——称量并赐给了我父亲们：——因为按我的时代\Emph{是}我的帮助——按我的疾病，我的药物——按我的美丽，我的装饰！\par

22. 我们\Emph{是}那扰乱了——传承和美好秩序的人——因为在温和的时代——看！我们却渴求严厉——好使您责备我们，仿佛\Emph{我们是}孩童一般！\par

\SourceRecord{Translated by J.T. Sarsfield Stopford.；Nicene and Post-Nicene Fathers, Second Series；Vol. 13.；Edited by Philip Schaff and Henry Wace.；Buffalo, NY: Christian Literature Publishing Co.,；1890.；<http://www.newadvent.org/fathers/3702c.htm>.}

% structure: p0001 source=h1 target=section reason=page-title

\section{尼西比圣歌}

% structure: p0002 source=h2 target=subsection reason=relative-heading-rank; base=h2

\subsection{圣歌 17.}

\Emph{关于\Place{尼西比}主教\Person{亚巴郎}}。\par

1. 主啊，请容许我这卑微的人也向祢的宝库投下我的分币，如同我们羊群中的商人，他善用了祢教导的塔冷通，增加所得，现已离去，进入祢的港湾。我要论及这位牧者，在已成为羊群之首的那位之下；他曾是那三位的门徒，现已成了我们的第四位导师。\Emph{R.，愿光荣归于祂，祂使他成为我们的安慰！}\par

2. 我要使他们在同一爱中发光，并将他们编作花冠，就是这位导师和他门徒的灿烂花朵与芬芳花朵，他如同\Person{厄里叟}留在他之后；因为他被选为角，受祝圣成为首领，被高举成为导师。\Emph{R.，愿光荣归于祂，祂使他成为首领！}\par

3. 他们在天上为羊群欢喜，因着他们牧养的那位牧者，羊群获得牧养；牧者的居所在他之下欢喜，因为他们看见他们等级的传承。祂取了他，将他作为心智安置在教会这个大身体中，他的肢体们围绕他，向他购买生命、教导和新饼。\Emph{R.，愿光荣归于祂，祂使他成为他们的宝库！}\par

4. 祂从众多牧者中拣选了他，因为他已证明自己的坚定；时间在羊群中考验了他，漫长的日子像坩埚一样证明了他；因为他以自身给出了证明，祂使他成为多人的墙壁。让你的斋戒成为我们国家的盔甲，你的祈祷成为我们城市的盾牌，你的香炉买回和好。\Emph{R.，愿光荣归于祂，祂圣化了你的祭献！}\par

5. 那位已与羊群分离的牧者，以属灵的牧场喂养他们，凭着他高举的牧杖，保护他们免受暗中的豺狼。填补你导师的位置，它渴慕他旋律的声音；将自己立为柱子，在颤抖的百姓城中；以你的祈祷支撑她。\Emph{R.，愿光荣归于祂，祂使你成为我们的柱子！}\par

6. 祂将手交给了他的门徒，将宝座交给配得上的人，将钥匙交给经过考验的忠信者，将羊群交给优异之人。按手礼属于你的手，赎罪属于你的祭献，安慰属于你的舌头。愿和平装饰你的主权；内在的看守和外在的会众。\Emph{R.，愿光荣归于祂，祂为喜庆拣选了你！}\par

7. 愿你的教导丰盛，多行少言！说很少的话，用劳动耕作我们的土地，这样因着多次耕作，稀少的种子可以变得丰盛，旧种子的增加，在我们中间有三十倍收成，你的新种子有六十倍。\Emph{R.，愿光荣归于祂，祂加倍百倍收成！}\par

8. 那曾反对你的愤怒止息了，因为平安完全涌向你；对你的嫉妒被浇灭，因为你的爱每时都在燃烧；你折断了嫉妒的弦，使它不能暗中击打任何人；混淆是非的诽谤，你的耳朵不转向它，因为公开的真理使你喜悦。\Emph{R.，愿光荣归于祂，祂装饰了你的肢体！}\par

9. 你要在百姓中提供建议，如同\Person{耶特洛}在希伯来人中；你要与那为你利益建议你的人完全同行，你要完全躲避那以其他方式建议你的人：\Person{勒哈贝罕}将成为你的鉴戒；你要选择有益的建议，你要拒绝嫉妒的建议。\Emph{R.，愿光荣归于祂，祂赐予安慰的建议！}\par

10. 那赐予你的恩赐，从高处飞降下来；你要用人名称呼它，你不要以另一种能力承受它，免得\Person{撒殚}以诡计来到它的位置，以为人子们将它给了你，致使这自由出生的恩赐为人的奴役服务。\Emph{R.，愿光荣归于祂，祂将祂的恩赐传了下来！}\par

11. 你的导师被描绘在你身上；看，祂的肖似完全在你身上。与我们分离的他与我们合一。在你内我们将看到那三位，那些与我们分离的卓越者。你将成为我们的墙，如同\Person{雅各伯}；充满温柔，如同\Person{巴布}；言语的宝库，如同\Person{瓦勒格什}。\Emph{R.，愿光荣归于祂，祂在一人身上描绘了他们！}\par

12. 我，羊群的渣滓，没有保留任何适当的事：我已用这两位的色彩描绘了这两位肖像，使羊群看见他们的装饰，羊群看见他们的美丽。而我，已变成一只有口才的羔羊，向祢，\Person{亚巴郎}的天主，以\Person{亚巴郎}的姿势赞美祢。\Emph{R.，愿光荣归于祂，祂使我成为祂的竖琴！}\par

% structure: p0016 source=h2 target=subsection reason=relative-heading-rank; base=h2

\subsection{圣歌 18.}

1. 啊，你，在你导师之后成为司铎，卓越者之后的杰出者，庄重者之后的贞洁者，节制者之后的警醒者。你的导师没有离开你；在生者身上我们看到逝者：因为看，祂的肖似被描绘在你内，他的足迹印在你身上，他的一切从你的一切中闪耀。\Emph{R.，愿光荣归于祂，祂将你赐给我们以代替祂！}\par

2. 果实，其中描绘了它的树，为根作证。直到现在，他的甘甜气味没有离开我们。你用身体行为展示他的话，因为你以行动成就了它们。在你的言谈中描绘了他的教义，在你的行为中是他的阐释，在你的成就中是他的诠释。\Emph{R.，愿光荣归于祂，祂使你的光辉卓越！}\par

3. 最后被高举的牧者，成为肢体的头，那获得长子权的年轻人，不像\Person{雅各伯}靠代价，不像\Person{亚郎}出于嫉妒，他的兄弟肋未人嫉妒他，而是像\Person{梅瑟}凭爱获得，尽管他比亚郎年长。在你的兄弟中，他们因你而欢喜如同因他。\Emph{R.，愿光荣归于祂，祂在同心合意中拣选了你！}\par

4. 在身体各肢体中没有嫉妒或妒忌，因为他们在爱中听从祂，他们被祂温柔地探访。头是肢体的了望塔，因为祂向四面观望。祂虽被高举，却因恩宠而谦卑，甚至俯就脚，为使他们免受害。\Emph{R.，愿光荣归于祂，祂将你的爱倾注在我们内！}\par

5. 这实在是小事，倘若背信被一位长者克服。老年以其明智屈服；青年在其时节得胜；因为一位年轻的斗士忍受了可憎的争斗，以充满背信的力量发动，如烟雾升腾又消散：其始即其终。\Emph{R.：愿祂受赞美，因祂向它吹气，它就消失了！}\par

6. 号角之声突然惊愕地召你出战。你上去如新\Person{达味}，被你制伏了第二个\Person{哥肋雅}。你并非未经争战，因为每日的隐秘争战，\Emph{你正在对抗\Person{恶者}}。隐秘的操练常在公开场合获得冠冕。\Emph{R.：愿祂受赞美，因祂拣选你作我们的荣耀！}\par

7. 面对试炼，\Person{约伯}操练他的身体和心灵，在诱惑中他得胜了。\Person{若瑟}在卧房中得胜；\Person{阿纳尼亚}及其同伴在火窑中，\Person{达尼尔}在狮窟中。\Person{撒殚}行事愚昧，因为在诱惑时，他公开证实了他们的胜利。\Emph{R.：愿祂受赞美，因祂加增他的羞辱！}\par

8. 那背信且急切的农夫，用左手撒播荆棘；正义的农夫对他满怀热忱，阻挡并砍断了他的左手。祂充满自己的右手，在心中播下生命之言；看哪！我们的悟性被祂的先知和祂的宗徒耕耘了。愿你耕耘我们的灵魂！\Emph{R.：愿祂受赞美，因祂拣选你作我们的农夫！}\par

9. 若你的话语太少，就以行为耕耘我们的地，使在多次耕耘中，茎和根得以坚固。一件美善的行为胜过听闻万言。你的种子将结实百倍，晚收六十倍，就连自生者也三十倍。\Emph{R.：愿祂受赞美，因祂加增你的收成！}\par

10. 光被昏暗是不恰当的，盐失去味道是不对的；头被玷污不相宜，镜子污秽也不相称。若药物失去效力，疾病也不能痊愈；若火炬熄灭，绊倒的也多。你的光将驱散我们的黑暗。\Emph{R.：愿祂受赞美，因祂使你作我们的灯！}\par

11. 为你任命经师和判官，督工和分配者，监督和官员：给每人分派他的工作，免得因挂虑而生锈，或因焦虑而分心，那心思和舌头，你用以献上恳求，为众百姓赎罪。\Emph{R.：愿祂受赞美，因祂使你的事奉显赫！}\par

12. 他洁净自己的心思，也洁净自己的舌头；他洁净自己的手，使全身发光；这对于司铎和其名号仍显不足，他奉献生活的身体。让他时时洁净自己；因为他作为中保，站立在天主与人类之间。\Emph{R.：愿祂受赞美，因祂洁净了祂的仆役！}\par

% structure: p0029 source=h2 target=subsection reason=relative-heading-rank; base=h2

\subsection{颂歌十九。}

1. 你回应\Person{亚巴郎}之名，因你被造为多人之父；但因为你没有配偶，如\Person{撒辣}之于\Person{亚巴郎}——看哪！你的羊群是你的配偶；以你的真理养育她的儿子；愿他们成为你的属灵儿女，并做应许之子，使他们成为\Place{伊甸}的继承者。\Emph{R.：愿祂受赞美，因祂在\Person{亚巴郎}身上预表了你！}\par

2. 贞洁的佳美果实，司铎职在你内得喜悦，在弟兄中最年幼，如\Person{叶瑟}之子；角满溢并傅抹了你，手降下并拣选了你，教会渴慕并爱了你；洁净的祭台为你的服务，伟大的宝座为你的荣耀，一切合一为你的冠冕。\Emph{R.：愿祂受赞美，因祂加增你的冠冕！}\par

3. 看哪！你的羊群，蒙福者，起来探望它，殷勤者！\Person{雅各伯}按次序排列羊群；排列这些会说话的羊，以纯洁启迪童贞少年，以贞洁启蒙童贞少女；兴起尊荣的司铎、温良的统治者、正义的百姓。\Emph{R.：愿祂受赞美，因祂以悟性充满你！}\par

4. 守护健康的羊，探望患病的，包扎破碎的，寻找迷失的；在圣经的草场上牧养它们，以教义之泉给它们饮水：让真理作你的墙垣，让十字架作你的杖，让诚实作你的平安。\Emph{R.：愿祂受赞美，因祂加增你的美德！}\par

5. 愿那与\Person{达味}同在的能力也与你在你的羊群中；因为若他从狮子口中夺回一只迷路的羔羊，你何其相配，崇高者，热心从\Person{恶者}手中夺取那些超越一切宝贵的灵魂，因为除\Person{基督}的血外，别无他物能赎买它们！\Emph{R.：愿祂受赞美，因祂被出卖却买赎了一切！}\par

6. \Person{若苏厄}服侍\Person{梅瑟}，因他服事的报酬，从他领受了右手。因为你服侍了一位杰出的长者，他也给了你右手。\Person{梅瑟}托付给\Person{若苏厄}一群羊，其中一半是狼；但交付给你的羊群，其中四分之一甚至三分之一是圣洁的。\Emph{R.：愿祂受赞美，因祂装饰了你的羊群！}\par

7. 愿\Person{梅瑟}的爱停留在你内，因为他的爱是明辨的爱，他的热忱是谨慎的热忱。当\Person{科辣黑}和\Person{达堂}分裂自己时，他将他脚下的地分裂；藉分裂使分裂止息。在\Person{厄耳达得}和\Person{默达得}身上，他显明他的善意完全在于全百姓都应说预言。\Emph{R.：愿祂受赞美，因祂在其善意中和好！}\par

8. \Person{厄里亚}的贫穷状况，\Person{厄里叟}爱之胜于财富；一个穷人给另一个穷人一件超越万有的礼物。因为你爱了你那暗中富足之主的贫穷，他言语的泉源要由你流出，使你成为圣神的竖琴，并内心里向他歌唱祂的善意。\Emph{R.：愿祂受赞美，因祂使你成为祂的宝藏！}\par

9. 无人嫉妒你的被选，因为你的领导是温柔的；无人因责备而发怒，因为你的话语播种和平；无人被你的声音惊吓，因为你的探访是愉悦的；无人因你的轭而叹息，因为它代替我们的颈项劳苦，减轻我们灵魂的重担。\Emph{R.：愿祂受赞美，因祂拣选你作我们的安息。}\par

10. 不可与强者争胜，亦不可对被遗弃者绝望；要柔和地教导富人，劝诫并争取穷人；对苛刻者要忍耐，对易怒者要宽容；藉善人捕获恶人，藉施舍者捕获掠夺者，藉圣洁者捕获污秽者。\Emph{R., 愿祂受讚颂，祂使你成为我们的猎者！}\par

11. 要带上一万种药剂，起来往病人中去；对患病者献上药物，对健康者献上预防；不可只献一种药剂，因为疾病或许不适宜：要献上许多疗法，好让疾病得医治；你也要从中学习经验。\Emph{R., 愿祂受讚颂，祂辛劳医治我们的创伤！}\par

12. 愿土地合乎你的心意，愿葡萄园合乎你的栽培；愿羊群居于你的住所，愿群羊在你的牧杖下平安！愿你成为伟大的元首，我们成为你冠冕上的宝石！愿我们在你内美丽，你也在我们内美丽！因为人民与司铎彼此相合时，二者皆美。\Emph{R., 愿祂受讚颂，祂在我们中播下了合一！}\par

13. 请听宗徒对他所许配的贞女所说的话：\Quote{我以嫉妒之情嫉妒你们，这是天主的嫉妒，不是肉性的，而是属神的。}\BibleRef{格后 11:2} 你们也当在纯洁中以这嫉妒去嫉妒，好让她知道她是什么，是属于谁的。愿她在你内珍惜，并在你内爱慕耶稣——真正的新郎。\Emph{R., 愿他的热忱是圣洁的那一位受讚颂！}\par

14. 主人如何，其习俗也如何：因为与懈怠的老师一起，她便是懈怠的；与崇高的老师一起，她便是卓越的。教会好比一面镜子，因为照着映照其中的面容，她便带上那面容的相似。正如君王如何，他的军队也如何；司铎如何，他的羊群也如何；按照这些，印记便打在它们身上。\Emph{R., 愿那按祂的肖像印刻她的那一位受讚颂！}\par

15. 没有遗嘱，他们离去了，那三位杰出的司铎；他们曾在盟约中默想，那两部天主的盟约。他们遗留给我们巨大的收益，即这贫乏的榜样。那些有福的、一无所有的人，使我们成为他们的产业；教会是他们的宝藏。\Emph{R., 愿那在他们内拥有其产业的那一位受讚颂！}\par

16. 司铎\Person{雅各伯}是尊贵的，她因他而尊贵如他：因为他将他的爱与他的嫉妒结合，披戴了敬畏与爱。与\Person{巴贝斯}，一位喜爱慷慨者，她用金钱赎回了俘虏。与\Person{瓦尔格什}，一位律法书记，她向圣经敞开了心。那么，愿藉着你，她的收益倍增！\Emph{R., 愿祂受讚颂，祂使她商贾众多！}\par

% structure: p0046 source=h2 target=subsection reason=relative-heading-rank; base=h2

\subsection{\Emph{Hymn 20.}}

1. 噢，成了新郎的童贞青年，请将你的心思稍微转向嫉妒，向着你青年之妻：斩断她在少女时代与众多他人所结的依附；斥责她，召集她的情感，好让她知道她是什么，是属于谁的。愿她在你内渴望，甚至爱慕基督——真正的良人。\Emph{R., 愿祂受讚颂，祂将她许配给祂的独生子！}\par

2. 噢，农夫，当嫉妒那莠子，它们已长出并纠缠在麦子中。拔除荆棘丛比拔除轻贱的野草更容易：若一阵轻风将它吹起，它便攻击庄稼并战胜它。那三位农夫所播种的，愿它三倍收回：三十倍、六十倍、一百倍！\Emph{R., 愿祂受讚颂，祂使你的收成丰盈！}\par

3. 一位新的牧者理当以新的方式照管羊群，并要知道它的数目，看出它的需要。这是一群用血买来的羊群，属于牧者之首。要呼唤每只羊，叫它的名字走过，因为这是一群名字被记录在生命册上的羊群。\Emph{R., 愿祂受讚颂，祂将索要牠们的数目！}\par

4. 看，你主的配偶与你同在！要保护她免受一切伤害，免受那些行事败坏、以自己的名号称呼会众的人。她丈夫的名字已印在她身上；她不可因别的名字而行淫，因为她不是以人的名字受洗；要以她受洗所用的名号宣认，即父及子及圣神之名。\Emph{R., 愿祂因祂的名号而受讚颂，她以这名号被呼召！}\par

5. 宗徒，她的许配者，曾为她嫉妒，免得她被名号败坏——但不是假名号，甚至连真名号也不可；甚至不可因\Person{刻法}的名号。那些真许配者将她的配偶的名号印在她身上；假许配者如同淫媒，将自己的名号印在羊群上。\Emph{R., 光荣归于你的名号，我们的造物主！}\par

6. 活物上的印记，我的弟兄们，无人公开毁坏；签署在信函上的名号，无人增添或更改：谁涂抹印记便是贼；谁更改名号便是伪造者。基督的名号已被更改；看，假名号已被印在已被败坏的会众身上。\Emph{R., 愿祂受讚颂，祂以祂的名号呼唤祂的羊群！}\par

7. 请看先知与宗徒，他们彼此多么相似！先知们将天主的名号印在天主的羊群上；宗徒们将基督的名号印在基督的教会身上。假许配者也彼此相似，因为会众以他们的名号被称呼，与他们行淫。\Emph{R., 愿祂受讚颂，我们因祂的名号而被圣化！}\par

% structure: p0054 source=h2 target=subsection reason=relative-heading-rank; base=h2

\subsection{\Emph{Hymn 21.}}

1. 若翰曾是火炬，揭露并斥责了放荡之人；他们急忙熄灭了火炬，好放纵他们欲望的情欲。要成为一盏明亮的灯，使黑暗的行为止息，好让你的教义升起时，无人敢在其光亮中，听从黑暗的情欲。\Emph{R., 愿祂受讚颂，祂使你成为我们的灯！}\par

2. 一个极大的祝福隐藏在其中，甚至是在厄里亚的责斥中。厄里叟服事了他，寻求他服事的双倍赏报。它赐给他双倍的光荣，因为他以双倍的程度披戴了他的德行。您既然爱了您富有的导师瓦尔格什的责斥，愿您继承他智慧的宝藏！\Emph{R., 愿赞美归於使您道理丰富的那位！}\par

3. 愿贪婪因您的斋戒而被克服，正如因达尼尔的斋戒一样！愿情欲在您的身体前蒙羞，正如在若瑟前蒙羞一样！愿贪财被您克服！正如在西默盎前被克服一样，愿您在地上捆绑，如同他一样，并在高处释放，效法他；因为您的信德如同他的信德！\Emph{R., 愿赞美归於那将祂的职务托付给您的那位！}\par

4. 愿您的贞洁如同厄里叟的，您的独身如同厄里亚的，您与眼睛立的盟约如同约伯的，您的仁慈如同达味的；愿您无嫉妒如同约纳堂，您的坚定如同耶肋米亚的，您的温良如同宗徒们！愿您拥有先知们的古物，宗徒们的新物。\Emph{R., 愿赞美归於用他们的宝藏充满您的那位！}\par

5. 愿您成为司祭职的冠冕，愿服务因您而闪耀！愿您成为长老们的弟兄，也是执事们的监督；愿您成为青年的导师，老年的拐杖和手；愿您成为贞女们的壁垒。愿盟约在您的交谈中得胜，教会因您的俊美而装饰。\Emph{R., 愿赞美归於那选您作司祭的那位！}\par

6. 愿在您的贫穷中，革哈齐家的可憎习俗归于无有；在您的圣德中，厄里家的可憎习俗被废除；在您的合一中，欺诈者依斯加略嘴唇的背信问候被除去！请倾泻我们的全部思想，从起初重新塑造它！\Emph{R., 愿赞美归於在您的坩埚中精炼我们的那一位！}\par

7. 愿在您的交谈中，那成为我们自由之主的\OriginalTerm{玛门}{Mammon}蒙羞！愿那惯于我们且令我们喜悦的疾病远离我们！请废除那维持了有害习俗的原因！恶事藉着习俗占据了我们的心：愿善事也藉着习俗占据我们！主，请成为我们援助的原因。\Emph{R., 愿赞美归於那为了我们的生命而拣选了您的那位！}\par

8. 愿恶习被断绝：愿教会不拥有财富；愿她拥有灵魂作为满足，若如此满足，将是奇妙无比的！愿她的亡者不被异教式地埋葬在失去希望中，穿着衣服、哀哭和悲叹；因为活人穿衣服，但亡者的一切就是棺材。\Emph{R., 愿赞美归於将我们再次归于尘土的那位！}\par

9. 邪恶的一个原因是情欲，也是厄里家的贪婪，革哈齐家的偷窃，以及纳巴耳的辱骂。请堵住这些可憎的泉源，免得有大水泛滥，从中流出污秽，甚至波及到您。愿主约束它们的泛滥！\Emph{R., 愿赞美归於使它们的泛滥干涸的那位！}\par

10. 因为对老人，将言语托付给他；对年轻人，命令他沉默；对来到您面前的外人，他从您的纪律中学习您，即谁先发言，谁第二和第三：若每个人守住自己的口，每个人知道自己的本分，他们将称您有福。\Emph{R., 愿我们的主成就您的愿望！}\par

11. 愿您真理的声音是单一的，您所采取的声音无数；真实性的形象在您心中，而在您的面容上是一切状态：悲伤、喜乐和软弱。对犯错者，显出您的愤怒；对贞洁者，显出您的喜乐。对天主性，保持单一；对人性，成为多面。\Emph{R., 愿赞美归於那与所有人成为一切的那位！}\par

12. 若您听到恶报，来自不说谎的诚实之人，请流出眼泪以扑灭在他人身上燃烧的火；愿智慧人与您一同祈祷，为有知识的人指定斋戒，愿您为那在罪中迷失的人居住在哀悼中，好使他藉着悔改回转。\Emph{R., 愿赞美归於那找到了迷失的羊的那位！}\par

13. 不要将您的耳朵给每个人，免得说谎者压倒您；不要将您的脚借给每个人，免得恶人误导您；不要将您的灵魂给每个人，免得傲慢者践踏您。使您的手远离虚假的人，免得他聚集荆棘到您手中。要深远又接近。\Emph{R., 愿赞美归於那虽远却近的那位！}\par

14. 看，新君的名声响彻世界！对被掠夺者是安慰，对掠夺者是恐怖。贪婪者要呕吐，好吐出他们吞下的一切。愿他们也在您面前恐惧，好使一位司祭和一位义王之间，旧习俗得以消除。\Emph{R., 愿赞美归於那发怒又转回施予仁慈的那位！}\par

15. 有人找着机会就冒险，有人强迫并强逼自己的意志。一人以为审判被保留，另一人以为根本不存在审判。有人偷窃并解渴，有人偷窃仍渴求偷窃。富人偷，穷人偷；但饥饿的人有分寸地偷，饱足的人毫无分寸地偷。\Emph{R., 愿赞美归於那查究一切意志的那位！}\par

16. 但现在祂给了机会，每个人都显明了他的意志，它是何种性质，像什么，以及他为自己选择了什么胜过什么。祂从每个人身上移除了诱惑，免得连那不恨祂的人也否认祂。祂给了我们机会，使我们明白，更好以为这能力\Emph{是}责罚，这责罚多多有益。\Emph{R., 愿赞美归於那为我们的益处而责斥我们的那位！}\par

17. 因为祂不愿以强制将祂的轭加在我们的颈项上；祂给了我们机会，我们便骄傲，好使我们既背叛又受惩罚，从而爱祂轻省的轭，选择祂宜人的杖。我们的安息对我们极其厌倦，因为在祂的强制中有安息，在祂的轭中有轻省。\Emph{R., 愿赞美归於那其劳苦是愉悦的那位！}\par

18. 整个世界如同一具躯体，已陷入重病之中；因为在异教的狂热中，它燃烧、憔悴并跌倒。慈悯的右手触摸了它，以怜悯对待它的灵魂；迅速切断了它的异教，因为那是它疾病的根源，它便被洁净、发汗并康复了。\Emph{R.，光荣归于治愈的手！}\par

19. 在你的日子里，大地将享有和平，因为它看见你充满和平。在你内，教会将被建造，披上它们的装饰，它们的书籍将在其中打开，它们的桌子将被摆设，它们的仆人将被装饰；从它们中要升起感恩，作为初果献给和平之主。\Emph{R.，愿祂复活我们的教会而受赞美！}\par

20. 愿你的祈祷升上天堂，愿和好也随之上升！愿上天之主将祂的祝福降在我们的[]上，将祂的安慰降在我们的苦难上，将祂的聚集降在我们的离散上；愿祂以祂的爱唤醒祂的嫉妒；愿祂的正义报复我们的耻辱，愿祂的恩宠抹去我们的不义！\Emph{R.，愿祂祝福祂的羊群而受赞美！}\par

21. \Emph{第一位}司铎和第一位君王，仿佛彼此描绘一般，如同天平般平衡。同样，瓦尔格什也是如此，那位君王的儿子也是如此，因为他们温良而平静。愿后者彼此相似；司铎成为光明，君王成为炽光，同样显赫的法官！\Emph{R.，愿祂光照了我们的灵魂而受赞美！}\par

22. 从君王职分出法律，从司铎职分出赎罪。两者皆温和可憎，两者皆强硬可悲。愿一者强硬，一者温和；在明智与审慎中，愿恐惧与怜悯相混合。愿我们的司铎职分温和，同样我们的君王刚强。\Emph{R.，愿祂调和了我们的援助而受赞美！}\par

23. 愿司铎为君王祈祷，使他们成为人类的屏障！胜利来自君王，信德来自司铎！愿胜利拯救\Emph{我们的}身体，信德拯救\Emph{我们的}灵魂！愿君王终结战争，司铎终结纷争！愿争辩与争吵止息！\Emph{R.，愿赐和平于万有者之子受赞美！为祢的恩赐，赞美归于祢！}\par

\SourceRecord{Translated by J.T. Sarsfield Stopford.；Nicene and Post-Nicene Fathers, Second Series；Vol. 13.；Edited by Philip Schaff and Henry Wace.；Buffalo, NY: Christian Literature Publishing Co.,；1890.；<http://www.newadvent.org/fathers/3702d.htm>.}

% structure: p0001 source=h1 target=section reason=page-title

\section{\WorkTitle{尼西比赞美诗}}

% structure: p0002 source=h2 target=subsection reason=relative-heading-rank; base=h2

\subsection{第35首}

\Emph{论吾主，并论死亡与撒殚。}\par

声音发出宣告：他们聚集前来，恶者的军队，连同他的差役。稗子的军旅完全聚集，因为他们看见耶稣已得胜，令所有左边的人忧伤，因为他们中没有不受折磨的。他们开始一个一个地述说他们所忍受的一切。罪与地狱战栗：死亡颤抖，死人反叛；撒殚因为罪人反叛他。\Emph{R.，愿光荣归于祢，因为恶者看见祢就惊慌！}\par

罪高声呼喊；她给她的儿子们——魔鬼和邪魔——出主意，对他们说：\Quote{军旅是你们行列的首领，那海已吞没了他和他的同伴；同样，我的儿子们，如果你们轻视这位耶稣，祂必消灭你们。你们曾用罗网擒获了撒罗满，因此你们受辱，竟被祂的门徒——渔夫和无知的人——胜过；因为看，他们捕获了曾被我们捕获的人。}\par

这是最大的，超越一切恶事（恶者论及我们的救主说）；因为祂夺去我们的一切还不够，还向我们开始了为约纳——阿米泰的儿子——报仇。因此，当祂擒获军旅并把他抛入海中时，正是为他报仇。约纳在三天后出来，升了上来；但军旅，即使在很长时期后也没有，因为海的深处奉命闭合在他身上。\par

在祂过去的经历之后，我用可口的饼试探祂，但祂不想要。我懊恼地努力学会一首圣咏，好藉祂的圣咏擒获祂：我暂停并第二次学会了它，但祂使我的第二次\Emph{试探}落空。我带祂上了一座山，将万有指给祂看；我把\Emph{这一切}给了祂，祂却不为所动。在亚当的日子，于我更好，他教导他并未给我多大麻烦。\par

恶者停止了活动，说道：\Quote{这耶稣是我懒惰的缘由；因为看！税吏和娼妓都投靠祂。我该为自己寻求什么工作？我原是万人的师傅，我该作谁的门徒？}罪又说：\Quote{我必须因此离弃并改变我的本性；因为这位来到的玛利亚之子，作为新的创造，创造了人类。}\par

饕餮的死亡哀叹说：\Quote{我学会了禁食，这是我素来不知道的；看！耶稣聚集群众，但至于我，在祂的宴席上，为我宣布了禁食。有一人堵住了我的口，我本是堵住许多人的口的。}地狱说：\Quote{我要控制我的贪婪；因此，饥饿是\Emph{我的}：这人如同在婚宴上得胜，当祂把水变成酒时，\Emph{同样}祂将死者的衣裳变为生命。}\par

此外，天主造了洪水，洗涤了大地，洁净了她的罪恶；又降下火与硫磺，为漂白她的污点。祂用火将索多玛人交给了我，用洪水交给了巨人。祂封住了散乃黑黎布军队的口，却张开了地狱的口。这些事以及类似的事，我都喜爱。\Emph{但现在}，代替致命的正义惩罚，祂在祂圣子内实行了因恩宠使死人复生。\par

恶者对同伴说：\Quote{先知和义人，我已见过；虽然他们的力量极其强大，但他们身上有属于我的气味；因为造人的材料与我们（的）天堂相近。这人却披上了亚当的身体，困扰我们，因为我们的酵母在祂身上无效。所以祂既是人，又是天主；因为祂的人性与祂的天主性互相融合。}\par

亚当被我看见，那涌出\Emph{人类}万族的泉源；他的儿女被我搜出，一个一个地考验。然而我从起初未曾见过一人，他的一部分出于天主，另一半是人。梅瑟，他光辉闪耀，我又试探了他，并使他在言语上犯错；但这人，不，祂的心思并非如此，因为祂思想的泉源极其纯洁。\par

肉体的贪欲存在于一切身体中；因为即使他们睡觉时，它也在他们内醒来。那在清醒时保持纯洁的人，我藉着梦境搅扰他。身体的渣滓在他内被搅动，通过内在隐秘的震动。睡着的和醒着的，我都同样搅扰。唯有这一位是独自保持纯洁的，我甚至在梦中也不能搅扰祂，祂在睡眠中也纯洁而神圣。\par

但祂的童年也与我见过的孩子们不同；因为我在祂身上没有看见\Emph{任何}属于我的部分。我害怕祂的童年；因此我煽动黑落德，使祂在婴孩中被杀。也因为祂逃脱了，我极其害怕，祂怎么发现了我们的奥秘！祂接受了贤士的礼物；祂藐视我们，离开并从我们的刀下逃脱了。\par

我见过孩童，义人的儿子；是的，也见过少年人，贞洁妇女的儿子；我从母胎中一个一个地撼动他们，看见他们内有我们的酵母。因为他们易怒、辱骂人，是的，也狂暴贪食；他们是通过教导才能成熟而变甜的果实。但这个人从祂\Emph{最初}的栽种，就是美好的果实，拥有甘甜，罪人藉此变得甘甜。\par

甚至在祂还是婴孩时，祂就因身上的光辉成为人类之师。连怀抱祂的司铎也惊异于祂。祂身穿老年人的智慧。若瑟远离祂：祂的母亲因祂的临在而荣耀。祂在童年时是每一个看见祂的人的帮助；从祂进入世界的那一天起，祂就使认识祂的人获益，祂以其美德成为人类的帮助者。\par

这玛利亚的果实，其葡萄所酿的酒不按本性，从我面前从哪里长出来？因为看！我处在疑惑之间。转身离开祂，我害怕，免得因祂的教导，那些因苦涩而获得的人变得甘甜。但再次践踏祂、压碎祂，对我是一个恐怖，免得祂反而成为罪人的新酒，当他们饮醉时，看！他们忘记了他们的偶像。\par

15. 看！我害怕这两件事，既有祂的死，也有祂的生。然后恶者的臣仆回答并劝告祂。虽然这两件事都令人痛苦，但对我们来说，麻烦稍轻一些的是，我们宁可选择祂的死，而非祂的生。让死亡告诉我们，是否曾有任何义人从起初就被唤醒？巨人的儿子们和有名望的人，没有一个能从她——吞食者的地狱中出来。\par

16. 人可以感知风的吹拂；但谁能看透玛利亚之子呢？因为当祂哭泣时，祂的眼泪夺走了我；而当我叫祂从圣殿跳下时，我以为祂是因恐惧而未跳下；然而当他们把祂从山顶扔下时，祂却腾空飞起。在井旁，当祂疲倦时，祂坐下。祂的变化我不明白，因为祂在旱地和水上行走。\par

17. 我看见祂饥饿，如同人子；但这被祂倍增的饼给消除了。从起初我试探祂，并来到祂那里；祂问我，好像不认识我；但这也被消除了，当祂显出祂知道我们的隐秘时。祂又选择了依斯加略，好像不认识他；然后转过来，显出认识他，尽管他正在捆绑和释放。我在祂身上错了，因为祂受了洗，又上来，淹没了我。\par

18. 但有一个记号我在祂身上看见，这远超一切使我振奋。因为当祂祈祷时，我看见祂，心里欢喜，因为祂变了颜色，感到恐惧；祂的汗珠如同血滴，因为祂感到祂的日子到了。这使我喜欢，远超一切，除非祂在那里欺骗了我。但如果祂哄骗了我，这对我和你们、我的臣仆，都一样。\par

22. 然后众魔鬼喊叫说：\Quote{我们在你身上看到的记号是可憎的，因为从起初从未这样临到你。你在敏捷的计谋中是卓越的；玛利亚之子夺取了我们的城市，而你却在延长你的话语。起来，出去，让我们与祂争战，因为这对我们是一种羞辱，我们这许多却被一个战胜。如果你疼痛或恐惧，就给我们战争的建议，留在后面。}\par

20. \Quote{这个耶稣，我要用祂自己的话教导祂，与祂争战；因为祂说撒殚是分裂的，自相纷争，就不能站立。虽然祂想与我们争战，祂却给了我们反对祂自己的武器：为我划分祂的门徒，因为如果你分裂他们，你就可以用这些征服他们，甚至用厄娃和蛇，这些软弱的势力，我用它们征服了第一个亚当。}\par

21. 死亡回答恶者，对他说：\Quote{你为何不照常等待？因为瞧，你所引诱的常是那些被轻视和微小的人；耶稣是超乎一切之上的，你用什么去引诱祂？祂的武器经验使你恐惧，就是祂受你试探时向你投掷的。你和你的随从，我们全军与祂——玛利亚之子争战都嫌不足。}\par

22. 那么我建议，如果这场争战容许我们做点什么：\Quote{进入那个门徒，你释放自己，让头与头说话；并释放你的全军，让它去煽动法利塞人。并且当心，不要照常争辩。如果你是神，就从这里下来，用爱亲吻他们，出卖祂；看，我们将把嫉妒和肋未人的刀剑加在祂身上。}\par

% structure: p0026 source=h2 target=subsection reason=relative-heading-rank; base=h2

\subsection{第三十六首颂歌。}

1. 我们的主抑制了祂的德能，约束了它，使祂活的死能给亚当生命。祂的双手被钉穿，代替那只摘取果实的手；在审判厅中祂被打脸，代替那在伊甸园中吃的手。又因他的脚从那里带走了亚当，祂的脚被刺透。我们的主被剥去衣服，为使我们谦逊；用苦胆和醋，祂使蛇的苦味变甜，这苦味祂曾倾注在人类身上。\newline{}（副歌）愿颂赞归于赐我胜利、使死者复起归于祂荣耀的那位！\par

2.（死亡说：）——\Quote{如果你是神，显示你的能力；如果你是人，就感受我们的能力。如果你寻找的是亚当，走开！因为他的过犯，他被关在这里；革鲁宾和色辣芬不能代替他偿还他的债。他们中没有可死的，能以生命代替他。谁能张开地狱的口，跳进去，把她所吞并紧握的他从她那里拉上来，且永远如此？}\par

3. \Quote{我是那征服了一切智者的人；看！他们在地狱中为我堆积在角落里。来，若瑟的儿子，进入看看可怕的事：巨人的肢体，参孙的强壮尸体，顽固的哥肋雅的骸骨；还有巨人儿子敖格，他为自己造了一张铁床躺在上面，我从那里把他推下扔下；那棵香柏树我放倒在阴间之门。}\par

4. \Quote{我一个人征服了众多，而有人想单枪匹马征服我。先知、司祭和名人，我都掠走了；我征服了军队中的君王、打猎中的勇士、德行中的义人。尸体的河流被我倾入地狱；虽然它们涌入，她仍感干渴。无论人近或远，终结都将他带至阴间之门。}\par

5. \Quote{我轻视富人手中的银子，他们的供物不能败坏我。奴隶的主人从未说服我，以奴隶代替他的主人，或以穷人代替富人，或以老人代替小孩。至于能驯服野兽的智者，他们的咒语不入我耳。所有人称我为憎恶劝说者；我只做奉命之事。}\par

6. \Quote{这是谁？祂是谁的儿子？祂的血统是什么，竟征服了我？家谱在我这里；看！我进去，阅读并研究了从亚当至今的名字，死者中没有一个我忘记的。家族接家族，看！它们写在我的肢体上。因为你，耶稣，我进去作了计算，好向你指明，无人逃脱我的手。}\par

7. 然而有两个人（我未说谎），他们的名字在地狱中我已忘记。因为哈诺客和厄里亚没有来到我这里。我在全世界寻找他们；是的，约纳曾下到那里，我也下去寻找，他们却不在。固然我猜想他们进入了乐园并逃脱了，但有一位强大的革鲁宾守护着它。雅各伯所见的梯子，或许他们正借此进入了天堂！\par

8. 谁能测量海沙，只洒落两粒？这收割中每天劳作的疾病如收割者，唯独我手持禾捆并收集；其他收割者匆忙劳作，却掉落禾捆。摘葡萄者忽略成串；\Emph{但是}两粒葡萄在我眼前逃脱，在那伟大的收成中，唯独我亲自采摘。\par

9. 我就是那在海上和陆地上追捕所有猎物的（死亡说）。空中的鹰隼来到我这里，深渊的巨龙也是如此：爬虫、飞鸟和牲畜；老人、青年和孩童。这些将说服祢，玛利亚之子，证明我的权能管治一切。祢的十字架如何能征服我？我早已藉一棵树得胜并征服了！\par

10. 但我还想再说，因为我不缺少言语；是的，我无需寻找言语，因为看！事实近在眼前召唤我。我不像祢向单纯的人许诺隐秘之事，说什么终将会有复活。如果祢果真全能，就赐予当下的质押，使祢遥远的应许也可得信。\par

11. 死亡结束了嘲弄的言论；我们主的声音响彻地狱，祂高声呼喊，将坟墓一一打破。颤抖抓住了死亡；那从未被照亮的地狱，忽然有荣光闪现，来自那进入其中的守望者，他们带出死者去迎接那已死却赐生命给众人的一位。死者出来，活人蒙羞，他们以为自己已战胜众人的生命赐予者。\par

12. 但谁给了我梅瑟的日子（死亡说）？他为我设了筵席。因为在埃及\Emph{被杀}的羔羊，赐给我每家中的首生子：成堆成堆的首生者，堆积在我地狱的门口。但这节日的羔羊却掠夺了地狱；祂夺取死者的名号，将他们从我这里带走。那羔羊为我填满坟墓；但这羔羊却倒空了满溢的坟墓。\par

13. 耶稣的死对我是一种折磨；我宁愿祂生而不愿祂死。这就是那死者，祂的死（看！）为我所憎恨；在所有\Emph{其他}人的死中我欢欣，但祂的死，就是祂的死，我厌恶；我盼望祂复活。祂活着时曾使三个死者复活并恢复生命；但现在藉祂的死，在我地狱门口，那些复活的死者——我正要关闭的——已践踏了我。\par

14. 我要赶紧关闭地狱之门，在这位死者——祂的死已掠夺了我——之前。凡听见的必会惊奇我的卑微，竟被一个外面的死者胜过。所有死者寻求出去，但这一个却竭力进入。生命的药剂进入了地狱，使它的死者恢复生命。是谁带入并向我隐藏了那活的火，使地狱寒冷幽暗的深处安息其中？\par

15. 死亡看见了地狱中的守望者；不朽的代替了可朽的；他说：混乱已进入我们的居所，因为在这两件事中有\Emph{我的}痛苦：死者出了地狱，而不死的守望者却进入其中。看！一个在这坟墓的枕边进入并坐于其旁，另一个同伴在祂脚边。我要恳求祂并劝服祂，带着祂的质押上去，往祂的王国去。\par

16. 仁慈的耶稣，请勿因我骄傲在祢面前所说的言语而发怒！谁见了祢的十字架，还会怀疑祢是人？谁见了祢的大能，还会不相信祢也是天主？看！藉这两件事，我已学会承认祢是人，同样也是天主！既然地狱中的死者不悔改，求祢上到生者中间，主，宣讲悔改。\par

17. 耶稣君王，请接纳我的恳求，并连同我的恳求收纳一份质押，就是那伟大的质押——亚当，在他内埋葬了所有死者；正如我接受他时，在他内隐藏了所有生者。我已将第一份质押赐予祢，就是亚当的身体；因此求祢上去统管万有；当我听见祢的号角时，我必亲手引领死者到祢面前。\par

18. 我们永活的君王已出去并上升，从地狱中，如得胜者。祸哉加于那些在左边的人；对于\Emph{邪恶}的灵和魔鬼，\Emph{祂是}忧伤，对于撒殚和死亡，\Emph{祂是}痛苦，对于罪和地狱，是哀哭。右边之人的喜乐今日已来到。因此在这伟大之日，我们要将这大光荣归于那死而\Emph{活}的一位，使祂能赐生命与复活给众人！\par

% structure: p0045 source=h2 target=subsection reason=relative-heading-rank; base=h2

\subsection{第三十七首赞歌}

1. 死亡为她，即为阴府哭泣，当他看见她的宝库已被倒空。他说：那么，是谁掠夺了你的财富？革哈齐偷窃而被发现；我每日偷窃，\Emph{但}偷窃尚未归咎于我。我被差遣到君王那里，在他们患病时，他们的卫士环绕他们，卫士也在他们的门口。君王的灵魂我攫取并离去。\Emph{R.，那以十字架折断死亡毒刺者，当受赞美！}\par

2. 所有妇女都为不孕忧愁；阴府却因自己的不孕而喜乐；若她生育，她便荒凉。全能的力量约束了她，即那\Emph{曾}不孕而冰冷的怀抱，它归还了债务，尽管惯于否认。黎贝加，当两个婴孩折磨她时，求死。那么阴府所受的痛苦何其大，当陌生的产痛袭击她时；死者被唤醒，冲出她的腹中。\par

3. 这难道就是我从\Person{依撒意亚}所听见的那句话？（但我曾轻视它）当他起来说：\BibleQuote{谁曾听说过这样的事？大地一日之内阵痛，在一小时内产生一个民族。} 这事真的发生了吗？还是为我们保留在将来？如果这就是，我自以为是君王乃是虚影；我不知道我只是在保管一笔存款。\par

4. 我从他，就是这位\Person{依撒意亚}，听到了两句不同的宣告。他说一位童贞女要怀孕生子；他又说大地要生产。但是看！童贞女生了祂，不生育的\OriginalTerm{阴府}{Sheol}也生了祂；两个违反本性的子宫，都被祂改变了；童贞女和阴府二者。童贞女因生产而喜乐；但阴府却因祂的复活而忧伤悲戚。\par

5. 我在山谷中看见了那位\Person{厄则克耳}，当他被询问时，他使死人复活；我看见成堆的骨头，它们就活动起来。\OriginalTerm{阴府}{Sheol}里骨头骚动，骨头寻找同伴，关节寻找配偶。那里没有询问者，也没有被询问者，问那些骨头是否活着。毫无疑问，\Person{耶稣}——万物的主宰——的声音使它们复活。\par

6. \OriginalTerm{阴府}{Sheol}看见他们心中忧伤，连忧伤的死者也欢喜起来。她为拉匝禄哭泣，当他出来时，她说：\Quote{平安去吧，\Emph{你}这活着的死人，被两家举哀者哀悼。} 内外都为他哀哭；因为他的姐妹们为他哭泣，当他进入坟墓到我这里时，他出来时我也为他哭泣。他死时，活人中有人哭泣；同样，阴府在他复活时也有极大的哀悼。\par

7. 现在我尝到了祂的忧伤的滋味，就是那为心爱者哭泣者的忧伤。那些被\OriginalTerm{阴府}{Sheol}所爱的死者，他们对他们的父亲是多么珍贵！我割下并带走的肢体，看！它们被剥夺并从我这里夺走。如果我因那青年的离去而受苦，他就是那重新获得生命的青年，有福的是那怜悯寡妇的主；在她的独生子中，祂赐平安给她那已荒凉的居所。\par

8. 看！我所加给人的、在他们心爱者身上的痛苦，最终将全然聚集到我身上。因为当死者离开\OriginalTerm{阴府}{Sheol}时，每个人都将复活，而只有我独自受折磨。那么谁能为我承担这一切，使我看见\OriginalTerm{阴府}{Sheol}被撇下孤独，因为那撕裂坟墓的声音，使她荒凉，释放她中间的死者？\par

9. 若一个人在先知书中阅读，他在那里听到正义的战争。但若一个人默想耶稣的故事，他学到的是恩宠和温柔的慈悲。若一个人认为耶稣是陌生的神，那是对我的侮辱。没有其他陌生的钥匙能配入\OriginalTerm{阴府}{Sheol}的门。只有创造者的钥匙，就是那已经打开、并且将在祂来临时打开它的钥匙。\par

10. 谁能将骨头连接起来，除了那创造它们的大能？什么能将身体的碎片重新结合，除了造物主的手？什么能恢复形体，除了创造者的手指？祂创造了、改变了、毁灭了，祂也能更新和复活。另一位神无法进入并恢复不属于祂的受造物。\par

11. 但若祂是另一位大能，我会很高兴祂来到我这里。祂要降入\OriginalTerm{阴府}{Sheol}的怀抱，并认识只有一位是天主。那些犯错并宣讲有许多神的凡人，看！他们在\OriginalTerm{阴府}{Sheol}被捆绑为我所得，而他们的神从未为他们忧伤。我只认识一位天主，我承认祂的先知和祂的宗徒。\par

% structure: p0057 source=h2 target=subsection reason=relative-heading-rank; base=h2

\subsection{第三十八首。}

1. 我的宝座在\OriginalTerm{阴府}{Sheol}中设立；一个死人起来，把我从上面推下。人人都惧怕我，我却不怕任何人。活人中是恐怖和烦恼，死人中是安息与平安。看！在一个被杀害的人身上，那俘掳阴府者进入了\OriginalTerm{阴府}{Sheol}。我曾俘掳所有人：我所俘掳的俘掳之子却俘掳了我。我所俘掳的那位，已经带领她离开，往\OriginalTerm{乐园}{Paradise}去了。\Emph{R., 祂十字架复活了阴府中的死者，祂是有福的！}\par

2. 众人都抱怨我；我只抱怨一个人。人间有谁像我这样正义？腐败触及了我的正直吗？我以慈爱对待所有人，恨我的人知道\Emph{这一点}；我一生不知什么是贿赂。我未曾接受君王的情面。借着我，平等被宣扬，因为在\OriginalTerm{阴府}{Sheol}中，奴仆与主人我使他们平等。\par

3. 我是在天主面前服务，祂不偏待人。有谁像我一样忍受，我做善事却被诅咒？我所作的善行被颠倒报答。虽然我的行为良好，我的名声却不佳。然而我的心灵安于正直：我在天主内安慰自己；因为虽然祂是善的，却每日被否认，并忍受\Emph{这一点}。\par

4. 我使老人脱离一切痛苦，青年脱离一切罪恶。我在\OriginalTerm{阴府}{Sheol}中平息隐秘的争论；在我们的土地上没有不义：只有\OriginalTerm{阴府}{Sheol}和天堂远离一切罪恶；这中间的大地，不义居住在其中。因此，明智的人要么上天堂，要么，如果那\Emph{太过}困难，就下到容易的阴府。\par

5. 为一个人因一个死者，每个人都急忙去安慰\Emph{他}。但为我，虽然我的许多死者已经复活，却没有人进来安慰我。\OriginalTerm{撒殚}{Satan}进来了，七祸已宣告反对他；虽然玛利亚之子大力践踏了他，\Emph{然而}他的灵高升；因为他是那虽受伤仍挣扎的蛇。我最好在这位以十字架征服我的\Person{耶稣}面前俯伏朝拜。\par

6. 当祂进入\OriginalTerm{阴府}{Sheol}的门，取代在祂来临前宣讲的若翰时，我将呼喊：看！那复活死者的来了；从今以后，我是祢的仆人，耶稣！因为身体我辱骂了祢，因为它遮掩了祢的\OriginalTerm{天主性}{Godhead}。不要生气，君王之子，对祢的宝库；按祢的命令我开了又关。虽然我的翅膀\Emph{极其}迅捷，但祢一点头，我就奔向各方。\par

7. 所有复活的人并非都是首生的，因为我们的主是阴府的首生者。一个死人怎能先于那使他复活的能力呢？有后来的成为首先的，有较小的成为首生的。因为默纳协虽是首生的，但厄弗辣因怎能获得长子的权利？如果次子被放在他前面，那么主和创造者岂不更要在祂的复活中超越\Emph{一切}！\par

8. 看，若翰作为先驱宣称他后来，虽然他生得较早；因为他说：\BibleQuote{有一个人在我以后来，却在我以前存在}。他怎能在他之前呢，即那他所宣讲的能力？因为一切因着\Emph{另一}事而成的事，都在\Emph{那另一}之后，即使它\Emph{看似}在先。因为使之\Emph{存在}的原因，在凡事上都\Emph{比它}更早、\Emph{更先}。\par

9. 亚当的原因比\Emph{所有}为他而造的受造物更早，因为创造者持续顾念他，甚至在创造之时。因此，虽然亚当尚未存在，他却比\Emph{所有}受造物更早。何况，我的主，这个祢的人性更早，它在祢的天主性中，从永远就与生祢者同在！愿赞美\Emph{归于}祢，并藉着祢归于祢的父，从我们\Emph{众人}！\par

10. 愿赞美\Emph{归于}祢，因为祢是首先的，在祢的天主性和祢的人性中！因为即使厄里亚是先升天的，他也不能阻止那为他而被接升的主。因为他的预像系于祢的真理：即使预像表面在祢的实现之前，实际上却在它们之后。受造物在亚当之前；他在它们之前，因为万物为他而造。\par

11. 我的主，请为我成就这个复活，不是由于祢的强制，而是出于祢的爱。因为祢的强制也使罪人得生命：依斯加略宁愿选择阴府的死亡，也不愿选择格赫纳的生命。请为我成就那出于祢仁慈的复活；即使祢的公义不允许，也请给祢的恩宠一个机会。只求祢为我记住，我在它里面寻求了庇护。\par

% structure: p0069 source=h2 target=subsection reason=relative-heading-rank; base=h2

\subsection{赞美诗第三十九首。}

1. 有赎者从圣人中来到我这里，但没有一个像玛利亚之子那样掠夺了我。因为看，厄里亚使死人复活；虽然他自身逃脱了我的手，但我仍有安慰，因为他所复活的死人，我从他那里夺走了。通过沙法特的儿子厄里叟，我仿佛被棍棒击打，因为他使两个死人复活。通过一根棍杖，我反而带走了先知和他所复活的死人。\Emph{应答：愿祂以声音劈开阴府坟墓者受赞美！}\par

2. 我害怕他，甚至革哈齐，当我看见他把棍杖放在少年身上。贼人取走了棍杖并返回；厄里叟来了，俯伏；像孩子一样低伏并升起，来回行走。我惊叹于我所看见的新奥秘，它只使一个少年恢复生命。\Emph{那时}我很好，当那些\Emph{只是}奥秘时，而不是现在，当死人反叛并征服了我。\par

3. 梅瑟，当我看见他脸上大能的光辉时，我害怕他；然而并未如我所怕的那样发生。他在阴府为我使尼散发芽；因为牧场，一个尸体的牧场，有六十万人倒下。——这个我所藐视的卑微者，治好了病人和患病者；祂给他人倍增饼食，但我们的饼\Emph{甚至}我们的饼，祂从我们口中夺走。\par

4. 在阴府曾有一场盛大宴席，当我吞下科辣黑和他的同伙时。撒殚为我带来极大快乐，当他在肋未人中制造纷争时。他在干旱之地为我涌出奶蜜之泉，当悖逆的会众下到阴府时。——看，义人活了并且出来了：梅瑟把活人送下去，但耶稣使死人复活并带上来。\par

5. 那时我很好，在那热忱者的日子，我在他们的刀剑中得喜悦。热忱者丕乃哈斯刺穿并给了我，在他矛尖上\Emph{作为}我的喜悦，齐默黎和苛次彼一同；在他枪尖上他把\Emph{他们}献给我。何时曾有两只肥牛犊，在矛尖上献祭？——但代替公主苛次彼，雅依洛的女儿被耶稣从我手中救出。\par

6. 亚郎的香炉使我恐惧，因为他站在死人与活人之间并征服了我。十字架更使我恐惧，它撕裂了阴府的坟墓。被钉在它上面的那一位，我杀了祂，如今我被祂所杀。被武装的战士战胜的人，他的羞辱并不很大。我的羞辱\Emph{是}比我的折磨更甚，因我被一个被钉的人所胜。\par

7. 丕乃哈斯的枪又使我恐惧，因为藉着他用枪所行的杀戮，他阻止了瘟疫。枪守护了生命树，它使我喜乐又使我悲伤；它阻止了亚当得生命，也阻止了死亡临于百姓。但刺透耶稣的枪，使我受苦；祂被刺，我呻吟。从祂那里流出了水和血；亚当洗净，活了，并返回了乐园。\par

8. 撒杜塞人仿佛是我的口，照我的心思与祂辩论，说根本没有死人复活。耶稣用一句话回答了他们，只有我明白；祂大声说出那可恶的话并使我忧愁：\BibleQuote{我是亚巴郎的天主，天主不是死人的天主。}那时我很好，这些\Emph{只是}话语，祂还\Emph{没有}向我显示死人真实的生命。\par

9. 农的儿子若苏厄杀了三十个王，为我填满了坟墓和坑穴；他毁灭了耶里哥，填满了阴府。但这位来的耶稣，荒废了\Emph{他们}死人的坟墓，并充满了上界各城。因此，看，他们名称相似，行为却不相似？\Emph{那}一位给了我阿苛尔的尸体，但\Emph{这}一位从我手中夺走了拉匝禄的尸体。\par

10. 梅瑟践踏了那埃及人，以他的温厚调和了正义。这新的法律从何而来，对我如此说：\Quote{若有人打你的面颊，你把另一面也转向他，并且注意你不可恨他。}\newline{}代替那充满熱心的强者，他曾践踏并杀人，一个慈悲的新人已为我们兴起。撒慕尔将阿加格斩成碎片，但耶稣却治愈了瘫痪的人。\BibleRef{玛 5:39}\BibleRef{玛 5:44}\par

11. 温柔慈悲，仿佛曾减少，看哪！在这时期已增长。而且它\Emph{当时}被憎恶，免得有人\Emph{藉着}它违犯诫命；因为撒乌耳和阿哈布因无慈悲而被杀，因为他们渴望怜悯恶人，而那些应受惩罚的人却没有被杀。在我的时代，耶稣改变了这事，赐生命给众人，并怜悯祂的杀害者。\par

12. 我记起三松那狮子的幼崽，他折断并给了我培肋舍特的柱子；还有那勇士阿贝乃尔，乃尔的儿子，为我擒获了那只敏捷的野羚羊，责鲁雅的儿子阿撒耳，并击打他，把他拋在地上。在圣殿中，贝纳雅公正地杀了雅各伯，正如所记载的。——因为正义已收敛了它的剑，从此悔改者要在恩宠中喜乐。\par

13. 达味用绳量了厄东人，一条线又一条线，毁灭了他们。那么，达味之子，祢何其慈悲！达味的正义是双倍的，当他杀死两行，却保留一行完全存活。——看哪！达味之子教导我们：\Quote{你要宽恕你的兄弟，直到七十个七次。}\newline{}在那里，正义被量度；但在这里，宽仁没有限度。\BibleRef{玛 18:22}\par

14. 达味满有热心和能力，他同时杀死了狮子和熊。他离开那大狮子，急忙去迎战那强壮的巨人。他用一块石子扑灭了他的光，他的灵魂离开他，他死亡了。但耶稣向那死去的青年呼喊：\Quote{青年！}\newline{}连死人在祂眼中都是睡着的。那青年祂救活了，并从我手中救出。祂为我淹死了那被轻视的猪群在海中。\BibleRef{路 7:14}\par

15. 肋未人因牛犊之故杀了他们的父亲和兄弟。依弗大亲手准备杀他的女儿。摩阿布王在城墙上献祭他的长子。面对他的刀剑，我喜乐。——但耶稣使刀剑变钝；是的，热病被斥责，那是阴府的姊妹：西满的岳母被治愈，但她被治愈的声誉使阴府痛苦。\par

16. 这耶稣，虽然祂是义者之子，祂所宣讲的全是恩宠。但对我来说，祂的恩宠是折磨。嫉妒是我们快乐的原因，因为嫉妒起初为我调和了第一次流血。为什么它在玛利亚之子眼中为有罪？祂来了，命令：\Quote{你不可对你的兄弟发怒。}\newline{}祂从兄弟间拿走了刀剑；而我从一开始就在加音的刀剑中得快乐。\BibleRef{玛 5:22}\par

17. 三松在骨中发现了蜂蜜——这岂不是一个奥秘？这耶稣为我们增多了奥秘。我落入了奥秘的波涛中，它们以比喻向我显示死者的生命，在各种奥秘和预像中。\Quote{食物来自吃者}是三松的比喻。但我却遭遇相反的事；因为吃者从食物中走出到我这里，因为从亚当中，看哪！\Emph{出来了}人子，祂毁灭了我。\BibleRef{民 14:14}\par

18. 义人们也多方掠夺了我，当他们宣讲死人复活时；但他们给我的悲哀掺入了极大的安慰。藉着阿撒和希则克雅的祈祷，我被喂养了死人，是的，我以尸体为宴席。厄里亚杀了巴耳的众先知，把他们交给我，那些靠依则贝尔的饼养肥的人。义人迫使我吞吃，但耶稣迫使我吐出我所吃的一切。\par

19. 我因所洒的血而害怕，梅瑟洒在每个门上的血；因为被杀者的血，却是拯救活人的血。古时我只怕那在门上的血，此外还有那在树上的血。被杀者的血是喜乐，如同馨香；但耶稣的血对我却是恐惧；因为我每次来闻祂的血，那藏在其中的生命之芳香令我战栗。\par

20. 司祭、大司祭、受傅者和君王，他们预表死人复活的预像，却从未藉着他们的十字架得胜。冠冕和冕旒\Emph{戴在他们头上}；当我与他们争战，我\Emph{有时}被打，\Emph{有时}也打人。但这个木匠的儿子，戴着荆棘冠冕，降卑并推倒我的骄傲，在祂的羞辱和死亡中：阴府见到祂，是的，从祂面前逃跑了。\par

21. 当大海见了梅瑟就逃跑，它因他的杖而害怕，也因他的荣耀。他的光辉、他的杖和他的能力，那被劈开的岩石也看见了。但阴府，当她的坟墓裂开时，她在祂——甚至耶稣——身上看见了什么？——代替光辉，祂穿上了死人的苍白，使她战栗。如果祂被杀时的苍白杀了她，那么当祂来复活死人时，祂在荣耀中，她怎能忍受呢？\par

% structure: p0091 source=h2 target=subsection reason=relative-heading-rank; base=h2

\subsection{赞美诗 40。}

1. 恶者察觉到他极大的羞辱，便在他的仆役面前自夸；他说大话来说服他们：\Quote{我所拥有的知识，少许出于本性；而大部分，是的全部，是习得而来。我为自己作了主人，并运用我的悟性。没有老师我就学会了一切；我以各样武器武装自己，并藉此赢得了我在人类中所渴望的冠冕。}\newline{}\Emph{副歌：愿颂赞归于那位前来解除罪恶罗网者！}\par

2. 在法利塞人中，我披上仇恨，为要与祂——玛利亚之子——争辩。愤怒如弓射出箭矢；狂妄辱骂祂；暴怒背叛祂；忘恩负义诽谤祂；嫉妒和妒火在\Emph{他们的}愤怒中与祂相争；亵渎拿起石头。医治者进来，站在病人中间，我煽动病人们与祂争辩。\par

3. 因为祂未受责难，我便以质问为避难之所。我屡次挑起事端，却见我的谎言被斥责，我的无耻被揭露，我的虚妄絮语被轻蔑。我投身于争辩的曲折之中。我无论何处与祂辩论，我的一切劳苦皆如糠秕，\OriginalTerm{真理之言}{word of truth}将其四散吹去。\par

4. 我看见一位战士，一位大能的主，以诡诈居于人内：【而\Emph{外面的}蛇使它恐惧。】他内心的情欲不断盘绕；他的嫉妒如蛇般嘶嘶作响。他生出致命的欲望，却又惧怕热病。命令如同药物，能平息那使人毁灭的嘲弄。唯有爱能折断那隐秘而苦毒的舌头之刺。\par

5. 谁比人更愚拙，竟为自己之外的事物操心？他每日检查他\Emph{箱中}的衣物，而蛀虫正潜伏在他肢体中。他修补衣服上的破口，灵魂却有了裂痕。他的房屋明亮，内心却黑暗。他关闭感官，却敞开窗户。他关上门，守护钱财；他的口张开，思想的宝藏却\Emph{被偷窃}。\par

6. 愚拙人看重他的牲畜胜过自己，因他关心财产胜过灵魂。他在田里撒好种，心中却撒稗子。他的理智被敞开推倒，却在葡萄园的篱笆上劳碌。他挑选并栽种葡萄树，而他的心思却是索多玛葡萄园的葡萄树。他赶走野驴不使其践踏庄稼，但林中的野猪却吞吃他的思想。\par

7. 我是人子的熔炉，他们的计谋在我里面受试炼。因此我得以编织欺骗。我教授\OriginalTerm{迦勒底之术}{Chaldean art}：因真事的应验，假事便被相信。在埃及中间，我遮蔽\Emph{人的}眼睛；我显现昆虫，\Emph{人}以为它们存在，其实并不存在。藉着遮蔽\Emph{人的}眼睛，我教授黄道十二宫的征兆，尽管它们并不在天上。\par

8. 因我的迅捷，我飞翔观看，预先告知占卜者；那些误解我的人视我为先知。但有时我胆大妄为，请求在一小时内向我启示隐秘之事，使真人藉我被证明，如同约伯，也使欺骗者如撒乌耳。对前者我揭露了他的巫术；对后者我炼净了他的真理，他便得称赞。\par

% structure: p0100 source=h2 target=subsection reason=relative-heading-rank; base=h2

\subsection{赞歌 41。}

1. 恶者说：我畏惧祂，甚至畏惧耶稣，恐怕祂毁灭我的技艺。看哪！我已存在数千年，从未有过安息。我未见有任何事物建立起来，使我转身离去。有一位来了，使不洁者纯洁：忧愁随至，因祂毁坏了我所建造的一切。我的劳苦和教训甚多，为要以诸恶覆盖一切受造物。\Emph{答：愿祂受赞颂，祂来揭露了\OriginalTerm{狡诈者}{the Crafty One}的诡计！}\par

2. 我与迅捷者比速度，且超越了他们：我发动战争；群众的骚动成了我的铠甲。我在百姓的骚动中欢喜，因它给了我现成的空间，因为群众的冲击是可怕的。藉着群众的力量，我筑起一座大山，一座高塔直达天上。若他们与至高者争战，何况\Emph{更加}战胜那在地上作战的祂呢？\par

3. 相机行事，待机而动，我谨慎地作战。百姓曾听说天主是唯一的，却为自己制造了众多神明。当他们看见\OriginalTerm{天主子}{Son of God}，便匆匆转向唯一的天主，以致仿佛承认天主，实则否认祂；仿佛出于热诚，实则逃避祂。如此，他们在一切悖逆之时，必被找到成为无神之人。\par

4. 看哪！我年岁久远，从未拒绝过任何婴孩。我常背负孩童的负担，为要从起初就使他们养成不良的习惯，使他们的过失随之成长。但有愚拙的父亲，不压碎我撒在他们儿子中的种子；也有\Emph{一些}人如同好农夫，从儿女的心思中拔除过失。\par

5. 我以懒惰如同锁链捆绑人，他们闲坐无事。我使他们的感官偏离一切善事：眼目远离阅读，口舌远离赞美歌唱，理智远离教义。对于有害而虚妄的寓言，他们\Emph{何等热衷}；对于空谈，\Emph{何等预备}！若\OriginalTerm{生命之言}{word of life}落在他们中间，他们要么推开它，要么起身离开祂的同在。\par

6. 人中间有多少\OriginalTerm{撒殚}{Satans}！唯独我，人人咒骂我。看哪！人的愤怒——\Emph{就是}每日磨碎他的魔鬼。魔鬼如同旅客，若被迫赶就会离去；但愤怒，尽管所有义人驱逐，却未从其处根除。众人憎恨一个软弱可怜的魔鬼，胜过有害的嫉妒。\par

7. 行邪术者与巫师一同蒙羞，他们每日驯服蛇。他里面的毒蛇却不在他能力之下；他驯服不了里面的情欲。隐秘的罪如角蝮，向他一吹，他便被灼伤。即使他凭狡计抓住毒蛇，错觉却暗中击打他。他以咒语安抚蛇；他以咒语激起对自己强烈的愤怒。\par

8. 我安置\Emph{我的}刺，坐下等待：谁\Emph{像我一样}对众人恒久忍耐？我坐在忍耐的人旁边，一步一步迷惑他，使他陷入绝望。那为过犯羞愧的人，习惯制服了他：我逐渐掌握他，直到他负轭，直到他进入其中习惯，甚至不愿出来。\par

9. 我看明，忍耐的人能制服一切。当我征服亚当时，他\Emph{只是}一人。我任凭他直到生养\Emph{儿女}，便为自己另寻任务，因为懒惰不合我口味。我数算海沙，为使我的心灵忍耐，并考验我的记忆是否足够，当人子增多之时。在他们增多之前，我已用许多事考验了他们。\par

10. 恶者的仆人与他争辩，并以他们的反驳驳斥了他的话。\Quote{但是，看！厄里叟使死人复活，在楼上战胜了死亡，并使寡妇的儿子复活。看！如今他在阴府中受束缚。} 但因恶者的推理十分有力，他用他们自己的话驳斥了他们的话。\Quote{厄里叟如何被战胜？看！在阴府中，他藉自己的骨头使死人复活。}\par

11. 如果厄里叟——他只有微小的\Emph{力量}——在阴府中却大有能力，且若果真在那里使一个死人复活，那么耶稣——那位大能者——的死亡在那里要唤醒多少死人呢！因此，即使从这一点也可以看出，耶稣远比我们这些同伴们伟大。因为，看！祂用祂的巧计欺骗了你们，而你们却不足以断定祂的伟大，竟将祂与先知们相比。\par

12. \Quote{你们的安慰是微小的\Emph{力量}，} 恶者对他同伙的人说。\Quote{因为那位使拉匝禄从死里复活的，死亡怎能敌得过祂？若死亡战胜了祂，那是祂愿意向死亡屈服；若祂果真愿意屈服，你们要大大害怕，因为祂不是白白死去。祂在我们心中种下了极大的恐惧，生怕祂临死时进入，唤醒亚当。}\par

13. 死亡从他的洞穴向外张望，看见我们的主被钉十字架，甚为惊奇，便说：\Quote{啊，唤醒死者的，你在哪里？你将成为我的食物，代替那甜蜜的拉匝禄，他的气味还在我口中。雅依洛的女儿要来观看你的十字架。寡妇的儿子注视着你。一棵树为我抓住了亚当；愿那为我抓住了达味之子的十字架受赞美！}\par

14. 死亡开口说：\Quote{玛利亚之子，难道你没有听见吗？梅瑟何等伟大卓越，超越众人？他成了神，行了神的事工？他杀了长子，救了长子？他使瘟疫离开生者？我与梅瑟一同上了山，愿荣耀归于祂的那位亲手把他交给了我。因为亚当之子无论变得多么伟大，他仍是尘土，必归于尘土，因为他出自泥土。}\par

15. 撒殚同他的仆人们前来，要看我们的主被投入阴府，好与他的谋士死亡一同欢喜；他看见主悲伤哀恸，因为死者因首生者的声音而复活，并从阴府走出。恶者起身安慰他的亲属死亡。你未能毁灭你所能毁灭的。正如耶稣在你中间，那已复活和活着的人将落入你手中。\par

16. \Quote{给我们打开，让我们看看祂，还要嘲笑祂：我们回答并说：\InnerQuote{你的能力在哪里？因为看！祂已经过了三天，让我们对祂说，啊，三日之主，你既使躺了四天的拉匝禄复活，就使自己复活吧。}} 死亡打开了阴府的门，我们主的面容的光辉从中照耀出来；他们像索多玛人一样被击打；他们摸索着寻找那已失落的阴府之门。\par

% structure: p0117 source=h2 target=subsection reason=relative-heading-rank; base=h2

\subsection{赞美诗第四十二首。}

1. 恶者哀号说：\Quote{如今哪里有我逃避义人的地方？我激动死亡去杀害宗徒们，好让我免受他们的打击。如今因他们的死，我反倒更受残酷的击打。我在印度杀害的宗徒，如今在厄得撒站在我面前：他完全在这里，也在那里。我去了那里，他在那里；这里和那里，我都找到他，并感到痛苦。} \Emph{应，愿那居于圣骨中的大能受赞美！}\par

2. 商人们携带的骨头，难道当时他们携带的是他本人吗？因为看！他们彼此得利。但对我，它们有何益处呢？是的，他们彼此得益，而我却从两者都受损害。但愿有人能让我看看依斯加略的钱袋，因为我藉它获得了力量！多默的钱袋却杀了我，因为住在其中的秘密力量折磨着我。\par

3. 蒙选的梅瑟带着骨头，怀着信心，如同带着财宝。如果伟大的先知相信骨头有益处，那么商人信得好，也自称为商人，信得好。那商人得了利益，变得强大，作了王。他的库房使我非常贫穷；他的库房在厄得撒开了，以益处使大城富足。\par

4. 我对这宝库感到惊奇，因为起初它的财富很小；虽然无人从中取走，它的财富之泉仍是贫乏的。但当众人围绕它，掠夺并带走它的财富时，它越被掠夺，财富就越增加。因为一个被堵住的泉源，若有人寻找它，当它被深深凿开时，便汹涌流出，更加丰富。\par

5. 显然，厄里叟在干渴的子民中是一道泉源；因为渴者不寻找他，他的流出就不大。但当纳阿曼寻找他时，他就丰富涌出，倾泻医治。那道泉源被带到另一道泉源中，他抓住他，把他浸入其中；因为在河中他洁净了癞病人。耶稣是恩惠的海洋，派那盲人到史罗亚，他的眼就开了。\par

6. 革哈齐拿着那曾使死人复活的杖，却不能使孩子复活。那著名的先知怎能被女巫召上来？我们是嘲笑撒乌耳的人，因为他问了一个魔鬼，却有两个魔鬼上来嘲笑他。从厄里叟的骨头也学习撒慕尔的\Emph{骨头}吧；因为虽然厄里叟的骨头使死人复活，巫师们却不能召上那活着而神圣的骨头。\par

7. 虽然我求了这个恩典，那位赐予万物者却没有赐给我。因为虽然邪魔被某些司铎、术士或巫师、加色丁人或占卜者的骨头所扰，但我意识到这只是嘲弄。我以两种方式使\Emph{人}犯错：要么我使宗徒们说谎，要么我使我的宗徒们如同宗徒们。\par

8. 魔鬼一党，看哪！已被掠掠；恶魔一党，受尽鞭笞：虽无人公开举杖，魔鬼却痛苦呼号；虽无人捆绑束缚，众灵却悬缚受刑。这沉默的审判，平静安稳，甚至不用审问，那唯一全能之力，看哪！住在第二位厄里叟的骨骸中。\par

9. 祂将审判权交给祂的十二位，使他们审判十二支派。如果他们要审判伟大亚巴郎的子孙，那么他们现在审判魔鬼也就不足为奇了。除非他们使钉十字架者完成将来的审判，否则他们将被我们的审判所证明。因为他们在宗徒——支派的审判者——面前呼喊得比我们更厉害。\par

10. 因为扫禄宗徒曾是豺狼，我以羊血喂养他；他长大强壮，成了独狼。但靠近大马士革时，狼突然变成了羊。他说宗徒将要审判天使；因为他以天使暗指司铎，正如经上所记。那么，如果他们如此强大，魔鬼就有祸了，因他们骨骸的击打！\par

\SourceRecord{Translated by J.T. Sarsfield Stopford.；Nicene and Post-Nicene Fathers, Second Series；Vol. 13.；Edited by Philip Schaff and Henry Wace.；Buffalo, NY: Christian Literature Publishing Co.,；1890.；<http://www.newadvent.org/fathers/3702e.htm>.}

% structure: p0001 source=h1 target=section reason=page-title

\section{尼西比斯赞美诗}

% structure: p0002 source=h2 target=subsection reason=relative-heading-rank; base=h2

\subsection{赞美诗第52首}

论撒殚与死亡。\par

1. 我听见死亡与撒殚在争论，在人中谁更有力量。\Emph{回应, 愿光荣归于祢, 万善牧之子, 祢从吞噬羊群的暗狼中拯救了祢的羊群, 恶者与死亡！}— 2. 死亡显示了他的力量，他征服了所有人；撒殚显示了他的诡诈，他使所有人犯罪。— 3. \Emph{死亡}, 恶者啊，除愿意者外无人听从你：无论愿意者还是不愿意者，都来到我这里。— 4. \Emph{撒殚}, 死亡啊，你的不过是暴政的力量：我的却是奸诈的网罗。— 5. \Emph{死亡}, 恶者啊，你听，凡狡诈者都折断你的轭：却无人能逃脱我的轭。— 6. \Emph{撒殚}, 死亡啊，你向患病者显示你的力量：但我对健康者却极其强大。— 7. \Emph{死亡}, 恶者不能胜过所有辱骂他的人：但凡是诅咒我的和咒骂我的，都落在我手中。— 8. \Emph{撒殚}, 死亡啊，你从天主那里得到了力量：我独自一人，无人帮助，当我引诱\Emph{人}犯罪时。— 9. \Emph{死亡}, 恶者啊，你像个弱者：而我则像君王行使我的统治。— 10. \Emph{撒殚}, 你是愚人，死亡啊，竟不知我何其伟大：我能捕捉自由意志，这至高权能。— 11. \Emph{死亡}, 恶者啊，你像贼，看！你四处游荡：我像狮子，击碎一切而无所畏惧。— 12. \Emph{撒殚}, 死亡啊，无人服侍或朝拜你：但列王向我献祭，如同朝拜天主。— 13. \Emph{死亡}, 许多人呼求死亡，如同呼唤一位仁慈的权能：却无人呼唤你，恶者。— 14. \Emph{撒殚}, 你不见此事吗，死亡啊，有多少人以各种方式呼唤我并献祭？— 15. \Emph{死亡}, 你的名字可憎，撒殚，你无法洗脱：你的名字人人诅咒，藏起你的羞辱吧。— 16. \Emph{撒殚}, 你的耳朵变钝了，死亡啊，以致听不见：全人类都在向你呻吟，你自己藏起来吧。— 17. \Emph{死亡}, 我向世界显露我的面容，因为我无诡诈：不像你，不能无诡诈而存留。— 18. \Emph{撒殚}, 你从未在任何事上超越我，因为确实：你和我一样，为世人所憎恶。— 19. \Emph{死亡}, 所有人都怕我如怕主人：至于你，他们恨你如恶者。— 20. \Emph{撒殚}, 死亡啊，他们恨你的名，也恨你的作为：他们恨我的名，却深爱我的快乐。— 21. \Emph{死亡}, 你的甘甜变成了牙齿的苦涩：灵魂的补赎始终紧附于你的欲望。— 22. \Emph{撒殚}, \OriginalTerm{阴府}{Sheol}被憎恨，因为她没有悔改：是吞没并封闭一切运动的深渊。— 23. \Emph{死亡}, 阴府是深渊，凡坠入者必再复活：罪被憎恨，因为它断绝了人的希望。— 24. \Emph{撒殚}, 我虽厌恶忏悔者，我仍\Emph{为悔改}留余地：你却断绝了死于自己罪中的罪人的希望。— 25. \Emph{死亡}, 起初，他的希望是因你而断绝的：因为凡你未引其犯罪者，死得安乐。— 26. \Emph{赞美祂}，祂使那两个被诅咒的仆人彼此对抗：好使我们看见他们，正如他们看见我们并嘲笑我们。— 27. 我们所见关于他们的，我的弟兄们，是一个保证：关于我们复活以后将要见到他们的情形。\par

% structure: p0005 source=h2 target=subsection reason=relative-heading-rank; base=h2

\subsection{赞美诗第53首}

1. 来吧，让我们听听他们如何争胜：那些有罪者从未征服过，也永不会征服。— 2. 死亡对恶者说：最终胜利属于我，因为死亡是终结的主宰，如同征服者。— 3. \Emph{撒殚}，这真是死亡，如果你能做到的话：通过情欲使活人致死。— 4. \Emph{死亡}，看！我注视死者，无论善与恶：那些轻视你的义人，却不轻视我。— 5. \Emph{撒殚}，身体的死亡只是暂时的睡眠：死亡啊，不要以为你就是死亡，你不过如影子。— 6. \Emph{死亡}，恶者啊，义人已征服了你，而且将要征服你：但这些征服你的人，看！我征服了他们。— 7. \Emph{撒殚}，即使你将义人致死，也不是出于你自己：是因为我所征服的亚当，他们才饮这杯。— 8. \Emph{死亡}，看！阴府充满了索多玛人、亚述人，还有洪水时的巨人，谁能像我？— 9. \Emph{撒殚}，死亡啊，这些全都由我杀死：是我使他们犯罪以致丧亡。— 10. \Emph{死亡}，撒殚啊，征服了你的\Person{若瑟}，我征服了他：在内室他征服了你，但我征服了他，将他投入坟墓。— 11. \Emph{撒殚}，死亡啊，以洒血征服了你的\Person{梅瑟}：他在\Place{埃及}征服了你，但在岩石旁，谁征服了他？— 12. \Emph{死亡}，撒殚啊，不畏惧你的\Person{厄里亚}：在\Person{依则贝耳}面前逃跑，因为他畏惧我。— 13. \Emph{撒殚}，死亡啊，以焚香抵御你的\Person{亚郎}：我给了他金耳环，他便铸了一只牛犊。— 14. \Emph{死亡}，撒殚啊，征服了你的\Person{约伯}，我征服了他：但在他征服你之后，我却征服了他。— 15. \Emph{撒殚}，以麻衣止住瘟疫的\Person{达味}：在屋顶上我征服了他，他曾征服\Person{哥肋雅}。— 16. \Emph{死亡}，摧毁巴耳庙宇、恶者殿堂的\Person{耶胡}：却不能摧毁我领域的堡垒阴府。— 17. \Emph{撒殚}，以他的判断从你口中夺回孩子的\Person{撒罗满}：在他老年，我使他成了偶像祭坛的建造者。— 18. \Emph{死亡}，撒殚啊，在金子方面蔑视你的\Person{撒慕尔}：我征服了他，这位征服者，他征服了贿赂。— 19. \Emph{撒殚}，死亡啊，在幼狮方面蔑视你的\Person{三松}：通过脆弱的器皿\Person{德里拉}，我使他套上了磨。— 20. \Emph{死亡}，恶者啊，\Person{约史雅}自幼年就轻视你：但即使到老年，他也未能抵挡我。— 21. \Emph{撒殚}，\Person{希则克雅}抵抗了你，死亡，当他克服生命的界限时：我误导了他，他忽略了奇迹，显出了他的宝藏。— 22. \Emph{死亡}，恶者啊，征服了你、赦罪并施洗的\Person{若翰}：我熄灭了那揭露你的火炬。— 23. \Emph{撒殚}，\Person{西满}征服了你，当他使那位有福的妇人复活时：在一个女人身上他征服了你，但通过一个女人我征服了他，使他否认。— 24. \Emph{撒殚}，宗徒和先知同声咒骂你，死亡：\BibleQuote{死亡，你的胜利在哪里？阴府，你的刺在哪里？}— 25. 被诅咒的仆人，你将你的主关在阴府里：天主憎恨你，人也同样，住口吧。— 26. \Emph{撒殚}，是那赋予一切生命者的旨意，将祂关在阴府：是你招致祂如此，当你使\Person{亚当}犯罪时。— 27. 哦，\Person{纳巴耳}在旷野辱骂其主的同伴：愿你的口受憎恶，你对祂说：\BibleQuote{俯伏朝拜我！}\par

% structure: p0007 source=h2 target=subsection reason=relative-heading-rank; base=h2

\subsection{歌 54。}

1. 听啊，自由，两位仆人的争论：他们如何彼此证实，本是无能。— 2. \Emph{响应，因祢的卑辱\Person{亚当}被高举，因祢的死亡他复活并重获\Place{伊甸}！}— 3. 如果恶者征服了你，耻辱甚大：死亡他的同伴证实了他，乃是软弱。— 4. 如果死亡再次制服了你，看！何等羞辱：因为恶者他的同伴嘲笑他，不过如影子。— 5. 他们的争论对你是镜子，在其中你可见：他们两者都如糠秕，在你气息之前。— 6. 是的，先知和宗徒，在他们的应许中：向你保证他们如花朵，在升起时凋谢。— 7. \Emph{撒殚}，死亡，你是他们憎恶的那位，生者和死者：因你解散并摧毁一切组合。— 8. \Emph{死亡}，撒殚啊，杀人的不是公开的死亡：你隐秘的死亡杀害人子。— 9. \Emph{撒殚}，我的名字不像你那样可憎，因为天使：在路上以撒殚的形象向\Person{巴郎}显现。— 10. \Emph{死亡}，撒殚啊，你的名字多么合适：你曾犯错并使不谨慎的\Person{亚当}偏离正路！— 11. \Emph{撒殚}，不要像无知者游荡，输了你的案子：死亡啊，如果你有能力就争辩回答。— 12. \Emph{死亡}，我知道你狡猾，撒殚：你能从沙中拧出圈套。— 13. \Emph{撒殚}，死亡，你的争辩结束了：因为被击败的人：当他的话枯竭结束时，开始咒骂。— 14. \Emph{死亡}，在万物中我是征服者，却被你击败？让\Person{亚当}说服你，他已被我征服，撒殚啊！— 15. \Emph{撒殚}，是我捆绑了\Person{亚当}，将他丢在你面前：那被我的诡计捆绑的壮士，你来制服了他。— 16. \Emph{死亡}，是我被重新加冕，在世上戴着冠冕：因为\Person{亚当}，壮士之首，我将他囚禁在阴府。— 17. \Emph{撒殚}，我以隐秘的死亡杀了他，就是\Person{亚当}当他犯罪时：死亡，你杀死了一个已死者，被我杀死的。— 18. \Emph{死亡}，恶者啊，在渴望征服中，你使自己被憎恨：因为你既是死亡又是撒殚，这对你\Emph{似乎是}小事。— 19. \Emph{撒殚}，死亡，你然后被沉默，如同弱者：因为无论在言语还是行为上，你都没有力量站立。— 20. \Emph{死亡}，恶者啊，你因你的恶得胜，如果你分辨，你的冠冕全然是羞辱，如果你察觉。— 21. 我将被击败，而你将被诅咒，撒殚：我宁可无知，也不为害。— 22. 愿正义者受赞美，祂分开了他们，尽管他们曾是同一心意：愿善者受赞美，祂使我们合一，当我们被分裂时。— 23. 我将通过祢的宽恕战胜恶者，全仁慈者：我将通过祢的复活战胜死亡，全生命给予者！\par

% structure: p0009 source=h2 target=subsection reason=relative-heading-rank; base=h2

\subsection{歌 55。}

1. 看哪！恶者责备死亡，又反被责备：从彼此，对着彼此，\Emph{是}他们的嘲弄。\newline{}2. \Emph{副歌：\Quote{愿祢受荣耀，万有之主之子，祢为众人而死：因为祢复活是为赐生命给众人，在祢降临之日！}}\newline{}3. \Emph{撒殚}：约纳征服了你，又从阴府返回，成了我的代言人，\Emph{在询问}为何罪人得免？\newline{}4. \Emph{死亡}：不要毁谤，啊，恶者，阿米泰之子：他显出忿怒的面容，好让他们更赞美你。\newline{}5. \Emph{撒殚}：你的一切劝说完全无力，啊，暴君死亡：因为你说的一切，没有一样令我喜悦。\newline{}6. \Emph{死亡}：真理之言何时曾令你喜悦？谎言者啊，深渊在你和诚实之间。\newline{}7. 我一生公义，无可悔改：我是那从你手中拯救世人之者。\newline{}8. \Emph{撒殚}：宣告你的悔改吧，死亡，你来得正好：看哪！连撒乌耳也在先知中，成了极大的笑柄。\newline{}9. 如果你，死亡，得以成义，那么对我自己：我也不断绝希望，同样悔改的指望。\newline{}10. \Emph{死亡}：我没有与我的主制造偶像，啊，憎恨你主者！看哪！你藉着死偶像杀害活人。\newline{}11. \Emph{撒殚}：死亡，我知道你是我的另一半，我也是你的一半：若我的一半悔改，它便悔改，但我惊异。\newline{}12. \Emph{死亡}：我是你分享中的同伴，但不是在罪中：被杀者是我的，杀人的是你的，是你使他们犯罪。\newline{}13. \Emph{撒殚}：我的狡猾为我而泣，当我与你争辩时：我的诡计为我哀恸，当我遇见你时。\newline{}14. \Emph{死亡}：行巫术的和占卜的，连同\Emph{他们的}一切过犯：你在世上点燃的火，我在阴府中已熄灭。\newline{}15. \Emph{撒殚}：你这滤出蠓虫、吞下义人的悔改者：贞洁者要撕裂你，他们从你腹中呼喊。\newline{}16. \Emph{死亡}：那是我保存所有义人的宝库\Emph{所在}：他们的复活威胁\Emph{祸害}给你，你曾迫害他们。\newline{}17. \Emph{撒殚}：贪婪者怀抱着一切受造物，在他腹中：看哪！他吐出来，说我被抢劫，失去我的财产。\newline{}18. \Emph{死亡}：在打击之前不要哀恸，因为还未\Emph{曾}临到\Emph{你}：日子将到，你要呼喊，我听见就欢喜。\newline{}19. 火要来，剥去你的皮：如同你用瓦片剥去约伯的皮。\newline{}20. \Emph{死亡}：懒惰的气味开始，仿佛要飘到我身上；那时我便间断，像一场梦，短暂时光。\newline{}21. 并非我无言可答，因此我沉默：我是为那荒废的时光悲伤。\newline{}22. 你话语\Emph{造成}的伤害极大：但愿我没有听见它！因为我全心专注于我的工作。\newline{}23. 这迷失的人类，因游荡的心思而败坏：懒惰和疏忽使他们负轭。\newline{}24. 欲望的疯狂竞标财富，并买下它：争执和夸耀是担保。\newline{}25. 我以坚韧为力量，进行我的战争：若我稍放松，我的统治便归无有。\newline{}26. 我以不断滴漏洁净岩石：因为不断滴漏能销蚀高山。\newline{}27. 习惯甚至能作自然的主宰：它训练并引领狮子，如同负重的牲畜。\newline{}28. 习惯、安息、增长，加上坚韧；由此自由被征服，尽管它极其顽固。\newline{}29. 若其意志坚定，它便挣断锁链；但若松懈，一张脆网也能捉住它。\newline{}30. 如果自由喊叫，我们便四散；但她若沉默，我们便聚集来嘲弄她。\newline{}31. 让我们停止多言，免得导致更多懒惰：同心协力攻打城墙，看哪！它倒塌了。\newline{}32. \Emph{撒殚}：你去照看疾病，我去照看网罗：因对我而言罪孽，对你而言瘟疫，是极大的慰藉。\newline{}33. 即使我停歇，我并未从忧患上停歇：因为我的意志从不休息，而是随时准备好。\par

% structure: p0011 source=h2 target=subsection reason=relative-heading-rank; base=h2

\subsection{圣诗五十六}

1. 自由与你搏斗，恶者啊：它能给你戴上口套，如果它愿意。—— 2. \Emph{R}，\Emph{光荣归于祢，因祢的胜利我们获得力量：因祢的复活我们甚至藐视死亡本身！}—— 3. 看！这两个又互相揭露，二者何等软弱：死亡提醒恶者你的大能（自由啊）。—— 4. 你的火在你巢中，死亡啊，你却不觉察：逝者的命运对你乃是倾覆。—— 5. 看！死亡和恶者宣扬你的大能（自由啊）：是的，恶者记起你的信德。—— 6. 如果这些反对你的成了你的同党：你的迫害者成了你的传报者，这是何等大事。—— 7. \Emph{D}，我承认，恶者啊，我像放高利贷般：积攒君王的财宝，直到祂的降临。—— 8. \Emph{S}，我，死亡啊，反而否认这属于天主：这奸诈的财宝，是我所储存的。—— 9. \Emph{D}，那么你的钱币是伪造的，撒殚啊：因为它不被收进天主的宝库。—— 10. \Emph{S}，我铸造新钱币，以君王之道：看！我的商人带来亏损，进入世界。—— 11. 天主从无中创造了万物：我从无中创造了巨大的罪。—— 12. \Emph{D}，愿你的口被封住捆绑，恶者啊，你这般大胆：看！竟将自己与创造者相比。—— 13. \Emph{S}，对我，死亡啊，我敢于冒险说话：你的舌头，连你的，也是奴隶，处于恐惧之下。—— 14. \Emph{D}，从此我们之间有深渊，撒殚啊：因为你疯狂地攻击你的主，看！—— 15. \Emph{S}，为何你怀疑，死亡啊，我们的和谐？做我们的同伙和肢体：看！我们为王。—— 16. 来，抽出我们的双剑，反对人类：我秘密地，你公开地，看！我们了结他们。—— 17. 罪和阴府也给了那二者劝告：说\Quote{如果你们分裂，就必败亡。}—— 18. 看水如何，如果分散，便流得低浅：但若汇聚，便获得力量，你们也是如此。—— 19. 如果分裂你们灭亡，如同软弱者：但联合一起你们为王，如同强者。—— 20. 爱熔化许多，如同在熔炉中：造成一个强大的整体，胜过一切。—— 21. 其中有智慧、狡诈、力量和能力：远比六十肘高的像更大。—— 22. 和好吧，让我们聚集起来，去反对那派：那派如果团结一致，就永不能战败。—— 23. 这些扰乱者谈论着，聚集并前来：主的日子，将把他们聚集，投入地狱。—— 24. 主啊，因祢的仁慈，当我复活时，我将朝拜祢：当我在祢的号声中受洁净时，我将赞美祢的圣子。\par

% structure: p0013 source=h2 target=subsection reason=relative-heading-rank; base=h2

\subsection{颂歌 57.}

1. 请听，我的弟兄们，听那死亡嘲笑邪恶者：因为它导致我们种族之首及其母亲犯了罪。—\newline{}2. \Emph{R.}，愿光荣归于祢，因祢的卑抑，撒殚被制伏：祢的贬抑高举了被贬抑的亚当。—\newline{}3. \Emph{D.}，你的大裸体将被亚当的子孙看见；正如你曾嘲笑他的裸体，当你使他犯罪时。—\newline{}4. 厄娃将离开那蛇，并斥责你：因为是你，龙，迷惑了她的单纯。—\newline{}5. 亚伯尔将看见他，即使加音也已来到你这里：他那愤怒的门徒将谴责他受诅咒的师傅。—\newline{}6. \Emph{S.}，诺厄战胜了洪水，如同死亡；我因含的口嘲笑他，当酒战胜他时。—\newline{}7. \Emph{D.}，诺厄并未受害，而是你的衣裳，你用它给他穿上：甚至诅咒，他穿上，成了奴隶。—\newline{}8. \Emph{S.}，罗特战胜了愤怒，这就是你的形象，死亡；我给我的女儿们如此建议，正如我所喜悦的。—\newline{}9. \Emph{D.}，罗特的妻子，她是你的器皿，听从了你的建议：愿你的一半干枯，正如你的整个器皿干枯了！—\newline{}10. 愿革亥拿倾覆在你头上：因为你的恶意倾覆了索多玛及其居民！—\newline{}11. 愿火的洪流在复活中向你涌起：因你向梅瑟和厄里亚煽动百姓！—\newline{}12. 愿义人最终嘲笑你，愿若瑟欢乐！他的兄弟们受你指使嘲笑他！—\newline{}13. 愿烟汽进入，窒息你的感官：正如海水窒息了恶人的感官！—\newline{}14. 愿贞洁的妇女也嘲笑你，因她们的建议：米德杨的女儿们嘲笑了愚昧的百姓！—\newline{}15. 愿火焰因三松的缘故在你头上点燃：因你借一个女人剃了他的发绺，那力量之狮！—\newline{}16. \Emph{S.}，撒乌耳，我以嫉妒征服了他，以巫术征服了你：因为他求问并带撒慕尔上来了，从他坟墓。—\newline{}17. \Emph{D.}，不要毁谤那活着的死人，因他并未上来：是你在那幻象中上来了，因为你是值得的。—\newline{}18. 愿诫命将你挂在火焰之上，邪恶者！因你使人将阿贝沙隆挂在一棵树上。—\newline{}19. 愿你在火中看见自己受卑抑，在卑贱妇女中间！因撒罗满因你被降卑，在亵渎的妇女中间。—\newline{}20. 愿公义按量还给你，像你点燃了它！因依则贝尔吞噬先知，你点燃了她。—\newline{}21. 愿你在火中公义地焚烧，因你使他们醉酒！厄里亚烧死了那两个人，当他们上去攻击\Emph{他}时。—\newline{}22. 愿煤也堆积在你身上！愿纳波特看见并欢乐，你在其中堆聚了石堆！—\newline{}23. 愿你在审判之日穿着嘲笑，在众目之前！因你借盗窃使革哈齐穿上癞病。—\newline{}24. 愿以闪电为标枪刺透你，撒殚！因你在约史雅的心中安置了\Emph{你的}标枪。—\newline{}25. 愿你在革亥拿的渣滓中沉没，撒殚！因你将耶肋米亚沉没在坑中的淤泥里。—\newline{}26. 达内尔逃脱了你将他投入的坑：愿他得安慰，永远看见你在火炉中！—\newline{}27. 愿你的恶行归在你头上，憎恨人类者！正如哈曼的恶行归在他头上，你的同伴！—\newline{}28. 愿王的净配嘲笑你，正如艾斯德尔：当你在审判日祈求她，为你求情时！—\newline{}29. 火释放了你所捆绑的义人：愿熊熊火焰成为你的大捆绑！—\newline{}30. 愿你被撕碎，愿七兄弟看见你的败亡：舍慕尼的儿子们，被你的狼群撕碎！—\newline{}31. 愿火焰战胜你的头顶，因你嘲笑：两个纳齐尔人的头，那不育妇人的儿子们！—\newline{}32. 愿火焰嘲笑你的头，因母亲和女儿：在若翰的头颅上得胜，当你激怒她们时！—\newline{}33. 火焰战胜了你的头，邪恶者：因你在诬告上战胜了若翰的头颅！\par

% structure: p0015 source=h2 target=subsection reason=relative-heading-rank; base=h2

\subsection{赞美诗 58.}

1. 看哪！死亡预先迅速行动，嘲笑撒旦：那个注定在最后成为笑柄的。—\newline{}2. \Emph{R.，光荣归于祢，祢藉祢的十字架刑征服了邪恶者，并藉祢的复活同样战胜了死亡！}—\newline{}3. 并为我们主的缘故，死亡向他说了诅咒：那人是祂蒙羞和十字架刑的原因。—\newline{}4. \Emph{D}．，愿火坑作你的坟墓，啊，撒旦：你亵渎了那来自坟墓的声音，它撕裂了坟墓。—\newline{}5. 我认识我的主，和我主之子，啊，你撒旦！你否认了你的主，并钉死了你的主之子。—\newline{}6. 这是适合你的名字，\Quote{杀害你主者}：当你所杀的那位显现时，祂必要杀你。—\newline{}7. 人人都要向你摇头，因为通过你，首领们向祂——生命之主——摇头。—\newline{}8. 你将是义人脚下一根压伤的芦苇：因为通过你，他们将一根芦苇放在那托住万有的祂手中。—\newline{}9. 祂被戴上荆棘冠冕，为要表明：祂拿取了达味家族的王国冠冕。—\newline{}10. 万王之王被戴上荆棘冠冕；但祂拿取了那羞辱祂之王的冠冕。—\newline{}11. 在他们给祂的戏弄袍子里，祂借此嘲笑了\Emph{他们}：因为祂拿取了司祭和君王的荣耀衣袍。—\newline{}12. 你的记忆相似醋，啊，你撒旦：你为生命泉的干渴献上了醋。—\newline{}13. 凡曾加强那击打祂的手的手，人人要举手反对你——万物藉祂的手而站立。—\newline{}14. 祂被手击打，而祂砍断了盖法的手：司祭职的手被砍断，在受傅的断绝中。—\newline{}15. 他们再次将祂绑在柱子上，如同为鞭打：那曾以云柱在前引导他们支派的祂。—\newline{}16. 柱子在柱子上，祂被鞭打：祂将自己从熙雍移走，于是它的倾覆来临。—\newline{}17. 当他们将两根木梁并在一起，形成十字架时：祂折断了它们，就是那两根杖，他们的守护者。—\newline{}18. 厄则克耳将两根木杖合在一处：在十字架的两根木梁中，他们的杖已止息。—\newline{}19. 那两根木杖，如同翅膀，承载了百姓：看哪！祂的两根杖都被折断，如同它的翅膀。—\newline{}20. 祂仁慈地张开了十字架的胸怀和翅膀：它的翼翎俯垂，承载万民，前往伊甸园。—\newline{}21. 它类似生命树，也类似其树干之子：它引领它所爱的，好让他们在其枝上饱食其果实。—\newline{}22. 去吧，邪恶者，哀号哭泣，为我也为你：因为我们中无一人能进入\Quote{生命园}。—\newline{}23. \Emph{S}．，既然你已承认，啊，死亡，来让我告诉你：你的这番言论，在我看来全是闲谈。—\newline{}24. 我要去查看我所设的网罗：你，死亡，也飞去照看所有病人。—\newline{}25. 我们的主使两者化为乌有，双方皆然：邪恶者将在此处被消减，死亡将在日后彼处被消减。\par

% structure: p0017 source=h2 target=subsection reason=relative-heading-rank; base=h2

\subsection{颂歌 59.}

1. 看哪！死亡为我们向撒旦施行报复：来听听他的羞辱并欢喜，因为他以我们的羞辱为喜。—\newline{}2. \Emph{R.，愿光荣归于祢，从祢的羊群，从祢：死亡和撒旦都屈服于祢的脚下！}—\newline{}3. \Emph{D}．，恶人将被正着悬挂，而你，头朝下：因为，相反地，你将西满钉在木头上。—\newline{}4. \Emph{S}．，关于其他一切，我保持沉默，死亡，因为我的时间渐逝：西满亲自嘱咐我，\Quote{这样钉死我。}—\newline{}5. 若是义人诅咒我，我还不至悲伤：死亡对我的诅咒，比地狱更糟。—\newline{}6. \Emph{D}．，我未提及我们主的羞辱，它太大，我的口难以言说：使我应权衡并比较祂的苦难与你的折磨。—\newline{}7. 祂将为祂的十二位设立十二个\Emph{审判}宝座：因为藉着十二支派，你，甚至你，将被定罪。—\newline{}8. 一条无价之绳将吊死你，啊，你撒旦：正如你的那位门徒以一条代价之绳自缢。—\newline{}9. 或许那边的地狱会仁慈地被清空：而你将在那里单独居住，与你的仆役们。—\newline{}10. 你的诅咒繁多，我怎能数算\Emph{它们}？看哪！你所有诅咒的总和，都在你的肢体上。—\newline{}11. 火中的恶人将刺伤你，你使他们变恶：他们将斥责你，\Quote{你为何带我们来这里？}—\newline{}12. 罪人将辱骂你，或许他们的威胁：对你比那边地狱的折磨更糟。—\newline{}13. 他们在那里对你都将是撒旦：正如你在这里对他们曾是那唯一的撒旦。—\newline{}14. 守望者将抓住你并抛下你，追忆：如何通过你，人们将他们的主从高处抛向深处。—\newline{}15. 所有人都将跑来用石头砸你，不忘：通过你，疯狂的人群跑来用石头砸他们的造主。—\newline{}16. 邪恶者，从众口将向你吐愤怒的唾沫：因为通过你，他们向祂吐唾沫，而祂的唾沫使盲人复明。—\newline{}17. 邪恶者，从众舌将向你发出所有诅咒：因为通过你，人们亵渎了祂，祂张开了哑巴的口。—\newline{}18. 有福\Emph{是}祂，祂为我们伸冤，虽默默无声，并激起死亡攻击邪恶者，使他落在其上！—\newline{}19. 我们高呼贺三纳，我的弟兄们，如基德红所\Emph{作}：\BibleRef{民 7:18}当他吹号时，压迫者彼此相攻！\par

% structure: p0019 source=h2 target=subsection reason=relative-heading-rank; base=h2

\subsection{颂歌 60.}

1. 哦，恶者突然遭逢何等惊愕，我的弟兄们：当那罪妇受责改正，获得智慧！—\newline{}— 2. \Emph{应：光荣归于那独自征服恶者的唯一者；也归于祂，是的，也归于祂，那战胜死亡的，予以宣认！}—\newline{}— 3. 恶者惊异道：\Quote{她的欢笑在哪里？她的香水在哪里？她的舞蹈、外表的装饰和内心的邪恶在哪里？}—\newline{}— 4. 她不再轻浮欢笑，转而流泪不止：她剪下头发，擦拭耶稣脚上的尘土。—\newline{}— 5. 任何教义在她身上都不持久，也不存留：她逃离了我们的教导，抛弃了我所教给她的一切。—\newline{}— 6. 她否认了我们和我们的相识，甚至仿佛从未见过我，将我的形象从她心中抹去。—\newline{}— 7. 耶稣的活酵飞向她，耶稣默不作声：但她却大胆闯入，虽然无人召唤她。—\newline{}— 8. 她忘记了我们多年的爱情，转瞬之间：她将之从我与她之间移除，置死亡于其间。—\newline{}— 9. 因为哭泣代替欢笑令她喜悦，泪如雨下代替脂粉，愁容代替装饰。—\newline{}— 10. 我使匝凯成为勒索者之首，使她成为荡妇之首；我的双翼，耶稣已折断。—\newline{}— 11. 倘若匝凯成为祂的门徒，倘若她成为祂的听众，从此他们束缚我的狡猾。—\newline{}— 12. 从此雕像成为笑柄，雕刻者成为讥讽对象，崇拜者成为笑料。—\newline{}— 13. 我遮蔽人的眼睛，使他们不能觉察那是雕像：耶稣开启他们的眼睛，使他们看出那是人手之作。—\newline{}— 14. 如果耶稣为自己拣选了宣讲者，那么我们的宣讲——那充满全世界的宣讲——就要归于沉寂。—\newline{}— 15. 看哪！加色丁人与占卜者，看哪！巫士与预言者一起被击打，司祭与所有恶者也被击打！—\newline{}— 16. 你们这些司祭已终结，从此断气，离开吧，占卜者！成为农夫，加色丁人同样要合上他们的书卷。—\newline{}— 17. 如果希伯来人成了祂的门徒——他们虽经历一切奇迹仍未被征服——那么万民中还有谁不服从祂呢？—\newline{}— 18. 如果祂开始矫正逆反之事，祂就使我们的言论归于虚无：从此祂——那斥责众人的——不会对我们犹豫。—\newline{}— 19. 既然我在所有庙宇中受敬拜，我们的耻辱就大于昔日的荣耀，因为众人唾弃我们的祭台。—\newline{}— 20. 祭献的肉变得可憎，被打碎：偶像被打碎，雕刻的像在锅下焚烧。—\newline{}— 21. 我们的一切工作成为笑柄和废墟：我们建造的一切，成为嘲弄；我们教导的一切，成为戏谑。—\newline{}— 22. 我辛劳教导他们的秘密奥义，即将被传扬在房顶上。—\newline{}— 23. 我以埃及人比任何民族都更自豪：因为他们甚至崇拜洋葱和大蒜。—\newline{}— 24. 看哪！我甚至害怕在这里，在迷惑如此深重的地方，真理会大大得胜，耶稣会在那里为王。—\newline{}— 25. 而且如果当祂还是婴儿时，逃往埃及，埃及人惊奇：甚至哄祂入睡——这个扼杀婴儿者，竟爱了他们的婴儿。—\newline{}— 26. 难道祂下去是为给她一个信物，如同订婚者：确保当祂成年时，祂也要娶她为妻吗？—\newline{}— 27. 法老不能站稳脚步，因为这不是一个口吃者，使他能欺骗祂，也不是一个奴仆，使他能向祂说谎。—\newline{}— 28. 梅瑟击打，埃及人反叛，他责罚百姓：希伯来人也反叛——耶稣被击打，却赐生命给众人。—\newline{}— 29. 这难以理解：祂不是用强力将祂的轭加在叛逆者身上；祂受责骂，却教导他人。—\newline{}— 30. 祂口中的唾沫，擦去并带走了亚当的羞辱：藉着击打祂的脸颊，祂从祂的门徒中根除了我们的怒气。—\newline{}— 31. 藉着祂所受的钉子，祂使我受苦。我因钉死祂而欢喜：却不知在祂的十字架上，祂正在钉死我。——\par

% structure: p0021 source=h2 target=subsection reason=relative-heading-rank; base=h2

\subsection{颂歌第六十一首}

1. 让我们以智慧倾听\Person{死神}，我的挚爱：他如何指责我们，为了我们的哭泣和哀悼。——\newline{}2. \Emph{R. 愿赞颂归于祢，祢降下来追随\Person{亚当}：寻见了\Person{亚当}，也在\Person{亚当}的子孙中。}——\newline{}3. 也许他正当地说：\Quote{你们无情地杀戮，看哪！你们却哭泣，好像有怜悯一样。}——\newline{}4. 你们使我成为一个残忍者，啊，你们这些杀人者：因为你们互相杀戮，无需我的帮助！——\newline{}5. 当\Person{死神}\Emph{只}渴望着来临，刀剑却已先他而至：让我们看看被杀戮者的血，在呼喊着反对谁。——\newline{}6. 被勒死的人，那些被窒息的人，在呼喊反对你们：因为那绳索，他们被勒死的绳索，使我感到羞耻。——\newline{}7. 他们甚至夺走了我的安息，因为若没有我：被勒死的人和被杀的人，怎能进入\Place{阴府}？——\newline{}8. 看哪！你们的婴儿被弃，如同在埃及的：你们把儿子祭献给魔鬼，啊，你们这些附魔者！——\newline{}9. 当\Person{死神}\Emph{只}渴望着尝一尝你们的尸体：\Person{加音}预先用人的血使我舒爽。——\newline{}10. 当我\Emph{只}渴望着耐心等待，直到\Person{亚当}应该死去：在我有权势之前，你们就给了我权柄，治理你们的身体。——\newline{}11. \Person{加音}用他的刀剑推倒了\Place{阴府}的门：因为它是关闭的，却在时候未到之前，他首先打开了它。——\newline{}12. 他借着踩踏，开辟了\Place{阴府}的道路，无需我的帮助：因为你们为我踏出了道路，看哪！我行走\Emph{在其中}。——\newline{}13. 九百年我坐着等待\Person{亚当}死去：但\Person{加音}连一天也不能容忍他的兄弟。——\newline{}14. 路上的强盗，比我还坏：我在打盹，而他们却在窥视杀戮。——\newline{}15. 看哪！你们的被杀者在坟墓里，你们的被谋杀者在路上；你们的被勒死在木桩上！——\newline{}16. \Quote{如果我背叛了我的主，并且杀了他：那么是谁杀了这些在这里的人？}\Person{耶胡}说。——\newline{}17. 如果我\Person{死神}已经取走了你们的逝者：那么那些被勒死的、被杀死的、被屠杀的，是谁杀了他们？——\newline{}18. 你们彼此互为\Person{撒殚}，\Person{恶者}被憎恶：你们彼此成为瘟疫，而\Person{死神}却被指责！——\newline{}19. 你们\Emph{自己的}意志对你们而言就是\Person{撒殚}，是的，也是杀人者：但所有人都抱怨\Person{死神}和\Person{撒殚}。——\newline{}20. \Person{死神}的毒药你们也互相给喝：看哪！除了我，你们还有多少死亡。——\newline{}21. 诡计、谋略、是的还有罗网、刀剑和毒药：从你们和你们里面，看哪！生出多少死亡。——\newline{}22. 审判厅里的法官，是第二个\Person{死神}：他为暗中的贿赂而杀人，我却不为任何东西。——\newline{}23. 我看见贿赂并对此惊奇，它跑得比我快：贿赂杀了多少人，却没有人察觉！——\newline{}24. 我羞愧自己如此笨拙地行事：如果我取走一具尸体，所有人都察觉了。——\newline{}25. 在房屋里哭泣，在街上哀号：甚至直到\Place{阴府}的门，他们都为我悲叹。——\newline{}26. 为你们自己悲叹吧，你们如此可憎，又恨我：\Place{阴府}从此要为你们悲叹，啊，杀人者！——\newline{}27. 用酷刑、鞭打和火，是的用石头砸：你们处死人的儿子，还自夸！——\newline{}28. 我比你们更谦逊、仁慈，也更敬畏：因为我怀着敬畏带走你们的逝者。——\newline{}29. 在床上我温柔对待病人：然后静静让他安睡，只\Emph{片刻}。\par

% structure: p0023 source=h2 target=subsection reason=relative-heading-rank; base=h2

\subsection{赞美诗 62.}

1. 看！死亡，沉默之王，抱怨，我的弟兄们：因为我们以绝望的哀号充满了它的居所。—— 2. \Emph{R., 愿极大的赞美归于祂，祂降临我们这里：受苦、复活，并在祂的身体内复活我们的身体！}—— 3. 当我们像疯子般在阴府门前哭泣时：请听死亡说话，它责备我们。—— 4. 死亡说，你们战胜了我，这使我羞愧：阴府的一半不足以容纳你们被杀的人。—— 5. 因为外邦人的尸体一起堆积在阴府：那里有两部分，死者和被杀者。—— 6. 然而我应当抱怨你们亏待了我，看！你们在哭泣：你们闯开了阴府的门，伤害了我。—— 7. 因为你们像一个婴儿，还在哭泣时：又笑起来，正如你们在死者之上。—— 8. 因为你们的哀悼没有分寸，也没有理解：在你们的笑声中——因为在我看来，你们就像一个断奶的婴孩。—— 9. 一时哭泣哀号，过一会儿：又戏笑放纵，如同孩童。—— 10. 因为你们不能成为完美的人：既不哭也不笑，如同谨慎的人。—— 11. 至于你们的书，我们感到忧伤，因为有人为它们劳苦：谁该为你们诵读它们，\Emph{甚至}神圣的圣经。—— 12. 诵读的人高声呼喊，因为你们是聋子：这\Emph{他们的}呼喊证明你们\Emph{如同}木石。—— 13. 因为既然诵读者和解释者在大声呼喊：你们的耳朵发沉，或是你们的心。—— 14. 因为如果\Emph{与你们同在}有一只耳朵，\Emph{敞开}于劝诫：就该少听多做。—— 15. 但因为\Emph{它}的听觉关闭，谁敲它：声音就返回到发出者。—— 16. 我并没有呼喊，我不聋：没有人为我诵读或解释，我不迟钝。—— 17. 从祂而来的气息命令我，\Emph{儿子们}，真理的天主：随着命令，也有完成。—— 18. 在我这里没有拖延，没有偏离：我知道甚至没有箭能超过我。—— 19. 但你们的声音被我蔑视，当你们在坟墓上哭泣：为你们离世的亲人，在绝望中。—— 20. 若可能或允许，当你们哭泣时：我要出来当面告诉\Emph{你们}。—— 21. \Quote{我努力交出死亡的账目：你们的声音打扰我，使我在计算中出错。}—— 22. 你们万民，不要让你们的理智变得幼稚：如同那从未有伟大智力的民族。—— 23. 在其中明智并不给予自己，如同愚人：因为他们的思想是黑暗，没有辨别。—— 24. 因为你们的婴孩和儿子，在复活时：他们将成为首先出来的，如初果—— 25. 然后\Emph{将出来}义人，如同得胜者：最后罪人将出来，如同蒙羞。—— 26. 因为虽然在眨眼之间他们得生命：却是按着次序，从阴府出来。—— 27. 先知和宗徒出来，以及\Emph{圣}祖：按着命令依次跟随。—— 28. 看！现在所撒的，是杂乱混合的：却以极大的次序归还，如园中的菜蔬。—— 29. 因为虽然在播种时，人可能混合各种种子：较早的苗妨碍后来的——30. 并且不像他们下去时混乱，上来时也是如此无序：从地里上来是确定的次序。—— 31. 看！我反对了我自己，我所说的：因为你们不明白的隐秘事，你们从我学到了。—— 32. 代替无益的眼泪，在坟墓旁：在\Emph{你们的}祈祷中，在教会中倾流它们。—— 33. 因为这些对死者有益，同样对生者：不要以折磨死者和活人的哭泣而哭泣！\par

% structure: p0025 source=h2 target=subsection reason=relative-heading-rank; base=h2

\subsection{赞歌 63.}

1. 谁能衡量亚巴郎的赏报？我惊奇他捆绑独生子的时候。—— 2. \Emph{R., 愿光荣归于你，使阴府中的死者复活的\OriginalTerm{圣言}{Voice}：他们已上来作祂圣子的宣讲者，祂使众人复活！}—— 3. 那时我急忙出来，要观看奇迹：他的刀如何抽出，对准他的爱子。—— 4. 我从四方聚集我众多的记忆：我收敛心神，惊奇那位显赫者。—— 5. 那么你们怎能诵读那伟大的故事？你们亲耳藐视了它的诵读。—— 6. 依弗大的剑责备那哀恸的人：他的女儿成为他从死得生的镜子。—— 7. 她为她父亲献出自己，所以你们当将你们的生命交托给万有之父，在\Emph{你们的}终局的希望中。—— 8. 那么，在母腹中你们不是经验了阴府的奥迹吗？然而在阴府你们比在母腹更得安息。—— 9. 你们反抗我强大的意志是顽固的：因为看！为了帮助它们，我夺走你们的离世者。—— 10. 摩阿布王亲手杀了他的儿子：那为离世者哀恸的人蒙羞。—— 11. 他是不敬神的人，看！按你们所读的：但你们是博士和教师，如你们所想。—— 12. 他忍受了，但你们在哀悼中暴怒：反抗万有的主的旨意，同时你们哭泣。—— 13. 然而我害怕放过约伯的故事：通过我这软弱的口，因为我不配。—— 14. 同样我避开\Emph{提及}他们的骨头：虽然我赞美祂，祂准许他们到我这里。—— 15. 不要因你们的罪侮辱你们的肢体：因为在阴府，恶人的骨头被鄙视。—— 16. 每当我看见恶人的身体：我践踏它并诅咒，甚至他的记忆。—— 17. 但每当我看见义人的骨头：我把它分开并尊敬它，向它敬拜。—— 18. 你们软弱的人不明白我所有的法令：在你们那里秩序混乱，因为你们是瞎子。—— 19. 我知道只有梅瑟像我一样尊重了那若瑟的骨头，我尊崇那若瑟。—— 20. 但梅瑟只对一个纯洁的身体行了这样的尊敬：但我对所有的义人的身体和骨头行尊敬。—— 21. 先知和宗徒的骨头明亮发光：所有的义人在黑暗中作我的灯。—— 22. 我敬拜那为我照亮阴府黑暗的祂：梅瑟的荣光如此伟大，对我如同太阳。\par

% structure: p0027 source=h2 target=subsection reason=relative-heading-rank; base=h2

\subsection{赞歌 64.}

1. 软弱的人啊，你们为何为你们的死者哭泣：他们在死亡中免于忧愁和罪过？—\newline{}2. \Emph{答：光荣归于那为众人忍受一切、为众人尝了死亡、为叫众人获得生命的那一位。}—\newline{}3. 我向你们揭示，连撒殚，虽然很满足：于你们的哭泣，却更嘲笑于你们的哀悼。—\newline{}4. 他戏弄地向我眨眼点头，如同一个小丑：\Quote{来吧，让我们嘲笑罪人，因为看啊！他们疯了。}—\newline{}5. 他们确已忘记了那为我为他们隐藏的火：看啊！这些愚人被哭泣灌醉，为他们的逝者。—\newline{}6. 不为没有准备而哭泣，好像我掠夺了并遣送了他们的死者，看啊！他们疯了。—\newline{}7. 恶人的灵魂将受磨难，直到审判日：而这些人像疯子一样在坟墓上哭泣。—\newline{}8. 他们不关心自己的罪过，恐怕明天：他们必以愧色，去加入他们的死者。—\newline{}9. 所有的人都将同样蒙羞，一家一家：在阴府中这些可怜人将无用地忏悔。—\newline{}10. 撇下醉汉和疯子，直到那一天：其中\Emph{每个}都要甩掉使他疯狂的酒。—\newline{}11. 我要去聚集他们，如同孩子：好让他们放浪疯癫，直到灭亡。—\newline{}12. 看！我已向你们揭示了奥秘，我同伴的秘密：因此出去，离开，改正，在悔改中。—\newline{}13. 离开我，我也将离去，我要处理我的事务：好使我开诚布公地向我的主交账。—\newline{}14. 我知道风吹走了我的话：因为你们正是我多次证明过的人。—\newline{}15. 我记起\Person{耶肋米亚}如何将顽固比作\OriginalTerm{印度人}{Indian}，他不改变\Emph{他的皮肤}，虽然它是自由的。—\newline{}16. 因为这也属于它，甚至自由：它以意志束缚自己，好像出于本性。—\newline{}17. 因为意志在自由人中如此有力：以致它通过其运作，可以被比作本性。\par

% structure: p0029 source=h2 target=subsection reason=relative-heading-rank; base=h2

\subsection{赞美诗第六十五首。}

1. \Emph{人}：死亡啊，不要轻视那亚当的肖像：它像种子被交付于土地，直到复活。—\newline{}2. \Emph{答：光荣归于你，你下降并潜入，追随亚当：将他从阴府深处拉出，带入伊甸！}—\newline{}3. \Emph{死亡}：我对这种子和你的话感到惊奇：因为看！五千年后，它仍未发芽。—\newline{}4. \Emph{人}：它目前的状态消逝，如冬天一样：而在复活中它像一把\Emph{谷物}进入生命的谷仓。—\newline{}5. \Emph{死}：我知道有收成季节，但我从未见过：死者被播种或收割。—\newline{}6. \Emph{人}：有一个收成要来临，死亡，它将使你赤露：\OriginalTerm{守望者}{Watchers}要出去如同收割者，使你荒凉。—\newline{}7. \Emph{死}：我何时成了农夫，而非葡萄园工？谁将阴府这压酒池变成了耕地？—\newline{}8. \Emph{人}：种子难道没有教导你，它腐朽而死：并断绝了希望，却从雨水重获希望？—\newline{}9. \Emph{死}：你们软弱的人看到了一个死里复活的梦：因为在醒着的时候，你们看不到复活。—\newline{}10. \Emph{人}：你的昏睡阻碍了你，使你看不到：许多高声呼喊的奥秘，关于复活。—\newline{}11. \Emph{死}：我知道种子活过来，但我从未见过：在阴府中生长的骨头，发芽并升起。—\newline{}12. \Emph{人}：你的全部言论就像你自己，因为看！\Person{厄则克耳}：已在山谷中教导你死者如何复活。—\newline{}13. \Emph{死}：我见过树木在夏天如何穿上它们的外衣：但骨头赤身裸体地被抛入阴府。—\newline{}14. \Emph{人}：\Person{梅瑟}以他的荣光击碎了你的心，死亡：亚当之子重获并穿上了亚当的荣耀。—\newline{}15. \Emph{死}：我们在阴府的法律是这样的：保持沉默：因为你们是言语，而我是行为，软弱的人啊。—\newline{}16. \Emph{人}：如果你是葡萄园工，老人如何被放过？那阻止你\Emph{夺取}他们生命的那一位，也使所有人活过来。—\newline{}17. \Emph{人}：子宫中的婴孩驳斥你，它\Emph{如同}被埋葬在那里：它向我宣告死里复活，向你宣告掠夺。—\newline{}18. 被藐视的花朵藐视你，因为它被闭锁并越过：虽然失丧，却不失丧，而是再次开花。—\newline{}19. 小鸡从蛋壳中啼叫，它被埋葬在其中：坟墓被一声巨响震裂，身体升起。—\newline{}20. 因为小鸡也是一具身体，在蛋壳里：看！它的身体向我们的身体宣告死里复活。—\newline{}21. 蝗虫推翻并了结了你的抗辩，死亡：因为它从尘土中出来教导死里复活。—\newline{}22. \Emph{死}：如果复活已经发生，我本可满足：因为复活的日子给我的困扰，不如你的审判。—\newline{}23. \Emph{死}：至高者的儿子是仁慈的，是的，良善而公义：不会严厉地报复我，为亚当的死亡。—\newline{}24. \Emph{死}：你难道没有理智理解这一点：你的父亲加给你这报应？\par

% structure: p0031 source=h2 target=subsection reason=relative-heading-rank; base=h2

\subsection{赞美诗第六十六首。}

1. 肃静，凡人（死亡说），片刻：并像我一样在阴府中如此沉默。— 2. \Emph{R., 光荣归于祢，守护者，祢降临到睡了的人之后，并从树木发声唤醒他们！}— 3. 你们悲伤，哀哭着为已离世的人，仿佛他来为我推阴府的磨。— 4. 我给疲倦之人的平安何其大：我不知疲倦，如你们一样，也不使他们疲倦。— 5. 我听到忘恩负义之人各式各样的咒骂：亚当的子孙如同亚当，对主忘恩负义。— 6. 你们的声音和行为彼此矛盾：用声音哀哭，却用行为天天争斗。— 7. 我听到哀哭，心想没有人劳作；我看到劳苦，心想没有人死亡。— 8. 人的挣扎使我以为他是不死的；他的嚎啕大哭使我以为明天他就不存在了。— 9. 且听，让我做你们的谋士，如果你们愿意：因为这两样，这些重担，非常苦涩。— 10. 稍微停下这劳苦和哀哭：要像凡人一样劳苦和哀哭，因为明天就消逝。— 11. 你们为逝者疯狂哀哭，为财产竭力劳苦。— 12. 死去的婴孩有福，他们是蒙福的：因为他们从你们所陷入的苦难中得了解脱。— 13. 让我去阴府，在那里说：\Quote{你们沉默的死者有福了，你们何等的安详！}— 14. 听我们话语的结论：如果有复活，就不要哭泣，也不要像外人一样劳苦。— 15. 你们像要永远活在这里一样挣扎，像永不会复活一样哭泣。— 16. 听我的话，如果你们有听的地方：预备你们的供给，以便我召唤时你们能回答。— 17. 因为我，我也听从那召唤我的，并将你们的身体连同你们的财宝恢复。— 18. 愿我们之间有平安，直到那一日：当你们出来时，我将呼喊说：\Quote{平安离去！}— 19. 来吧，你我现在就要将光荣归于那掌管死亡和生命的主，愿祂赐予帮助。— 20. 愿我们从众人而来的赞美归于祢，主，永生的祭献！祢藉着祢身体的祭献，赐予活人和死人生命。— 21. 赞美那穿上我们身体、死而复活的主：祂在我们内死了，我们在祂内活着，愿派遣祂的那位受赞美！\par

% structure: p0033 source=h2 target=subsection reason=relative-heading-rank; base=h2

\subsection{赞美诗67}

1. 来，让我们听死亡如何定罪百姓：他们的刀剑比死亡更残忍，攻击义人。— 2. \Emph{R., 光荣归于祢，祢藉着祢的祭献，赎回了我们的耻辱：祢的死代替了所有的死亡，使祢能复活众人！}— 3. 的确不是死亡钉死了耶稣，而是百姓：百姓多么可憎，甚至比我还可憎！— 4. 他们将耶肋米亚投入坑中，泥泞的坑：但我在阴府却将尊荣分配给祂的骨骸。— 5. 他们用石头将纳波特砸死，如同狗一样：我多么仁慈，连狗也从未扔过石头！— 6. 希伯来妇女在饥荒中吃掉自己的孩子：阴府却是仁慈的，她轻易地交出并归还他们。— 7. 藉着厄里亚的手，我将儿子还给寡妇；藉着厄里叟的手，将心爱的还给叔能妇女。— 8. 希伯来妇女出于贪婪吃掉孩子：阴府交出死者，学会了谨守斋戒。— 9. 阴府并非真正的阴府，不过是其仿冒品：依则贝耳才是真正的阴府，她吞噬了义人。— 10. 先知们的儿子和先知，她杀害并推倒；厄里亚逃到天上，躲避她的狂怒。— 11. 百姓之中有多少死亡代替了一个死亡！又有多少阴府代替了一个！— 12. 在以色列，撒玛黎雅和依次勒耳是她的女儿；在犹太，熙雍和耶路撒冷是她的姊妹。— 13. 在犹太和以色列的先知和义人，都淹死在这两个深渊中。— 14. 那么为何阴府被憎恨，唯独她一人被恨？虽然有许多比她更可憎的？— 15. 犹大之人的死者，对我极为可憎：是的，他们的骨骸在阴府中被我厌恶。— 16. 但愿我有办法将他们赶出去，将他们的骨骸从阴府赶出，因为他们使她腐烂。— 17. 我惊异于圣神，祂竟如此居住在其中的一个百姓，他们的气味如同他们的言行一般恶臭。— 18. 葱和大蒜是他们行为的宣告：正如这污秽之人的食物，他们的悟性也是如此。— 19. 藉着所有俯伏敬拜祢父的人之祈祷：求祢怜悯祢的敬拜者，他对祢的爱忘恩负义。— 20. 来自希伯来人和阿兰人，也来自守护者：愿赞美归于祢，并藉着祢归于祢的父，也愿光荣归于祢！— 21. 因为我有一张口面对那无口的死亡：愿那作为一切口的圣子，使我的过犯不触犯祂的父！\par

% structure: p0035 source=h2 target=subsection reason=relative-heading-rank; base=h2

\subsection{赞美诗68}

1. \Emph{人}\Quote{：啊，死亡，不要因义人而夸耀：你主的众子，奉祂的命令，来与你\Emph{同住}。}——\newline{}2. \Emph{众}\Quote{：愿光荣归于祢，因着祢的命令，死亡曾掌权；因着祢的复活，它已被贬抑卑微！}——\newline{}3. \Emph{死亡}\Quote{：在此我便极其伟大，如你所说：我虽是仆人，却践踏自由的人。}——\newline{}4. \Emph{人}\Quote{：亚当本被拣选为统治者；在他轭下，死亡和你的同伴恶者成了仆人。}——\newline{}5. \Emph{死亡}\Quote{：这正是我们的骄傲：看！奴隶成了主人：死亡和它的同伴撒殚践踏了亚当。}——\newline{}6. \Emph{人}\Quote{：看！你和你的同伴，受诅咒的仆人，所受的贬抑！哈诺客如何践踏你们俩，上升高处而掌权。}——\newline{}7. \Emph{死亡}\Quote{：虽然哈诺客使我忧伤，但我仍有安慰：因为在阴府中诺厄的尘土上，看！我践踏。}——\newline{}8. \Emph{人}\Quote{：死亡，在人类面前战栗吧！因为虽为仆人，他统治的轭却管辖一切受造物。}——\newline{}9. \Emph{死亡}\Quote{：那么我欢喜，因为被我战胜的并非卑劣的\Emph{敌人}：因为按被征服者的伟大，征服者也伟大。}——\newline{}10. \Emph{人}\Quote{：死亡，你的声音唱得胜利，对着义人：因为哈诺客和厄里亚已折断你的双翼。}——\newline{}11. \Emph{死亡}\Quote{：我知道\Emph{如何}用安慰来衡量悲伤：代替两个，看！许多已来，且正来临。}——\newline{}12. \Emph{人}\Quote{：所有已来和将来到你那里的人，如同旅客寄居，又像拉匝禄离开你的居所。}——\newline{}13. \Emph{死亡}\Quote{：你的话不伤我，反医治我：因为反抗我的拉匝禄，我又降服了他。}——\newline{}14. \Emph{人}\Quote{：死亡，请回答并辩论，是什么约束了他，使他被复活，除非是一个奥迹，预示他的复活。}——\newline{}15. \Emph{死亡}\Quote{：你以争辩闻名如同闲人，而我却辛劳于分辨和实行我的任务。}——\newline{}16. \Emph{人}\Quote{：你早已准备好辩论，是什么阻止了你？我们复活的真理以其名声约束了你。}——\newline{}17. \Emph{死亡}\Quote{：你使我被你们憎恨，虽然我并不可恨：我是那使你们年老者和忧苦者得安息的一位。}——\newline{}18. \Emph{死亡}\Quote{：你使我如同惹麻烦者，啊，你们必死的人：亚当将死亡带给你们，而我承受责备。}——\newline{}19. \Emph{死亡}\Quote{：我将温和地揭露你们，因为我是奴隶，而你们正是那些因自己的罪使我成王的人。}——\newline{}20. \Emph{死亡}\Quote{：亚当的意志激起了我，因我本在安息：我本是死的，你们却使我活跃，好叫你们因我而死。}——\newline{}21. \Emph{死亡}\Quote{：我控告那些说谎者，他们杀了人却否认：因为亚当杀了自己，却归咎于我。}——\newline{}22. \Emph{死亡}\Quote{：争端的开端是那受诅咒的蛇，它已正当地被击伤：它爬行、进入，并在你我之间设立了仇恨。}——\newline{}23. \Emph{死亡}\Quote{：撒殚过去了，你们却是针对我被激起：去吧，与那使你们悖逆的恶者争斗。}——\newline{}24. \Emph{死亡}\Quote{：他是我的同伴，我并不否认\Emph{它}，但尽管他极受憎恨，何必因他而责备我。从今以后，我否认他。}——\newline{}25. \Emph{死亡}\Quote{：请听我的话，啊，必死的人，我要安慰你们：我曾使你们受苦，我承认从死者中来的生命。}——\newline{}26. \Emph{死亡}\Quote{：因为开始有预备的声音潜入我的耳中：那预备好吹响的号角之声。}——\newline{}27. \Emph{死亡}\Quote{：听我的话，将许多油放入你们的灯中：因为从我这一边，对你们没有阻碍。}——\newline{}28. \Emph{然而}\Quote{：要知道，即使我说了这些，你们的声音在阴府的孤寂中仍是悦耳的。}——\newline{}29. \Emph{死亡}\Quote{：因为人已由我衡量，他的平安是大的：因为蛇、鱼和鸟都来迎接他。}——\newline{}30. \Emph{死亡}\Quote{：但奇妙的是，对守望者们，他的交谈也是宝贵的：是的，在革黎纳的恶者，也渴望他的同在。}——\newline{}31. \Emph{死亡}\Quote{：你们将从死者中获得生命，啊，必死的人，而我——这被剥夺的——将在阴府中被剥夺。}——\newline{}32. \Emph{人}\Quote{：愿赞美从众人上升于祢——祢使一切活，并从各处聚集亚当的尘土！}\par

\SourceRecord{Translated by J.T. Sarsfield Stopford.；Nicene and Post-Nicene Fathers, Second Series；Vol. 13.；Edited by Philip Schaff and Henry Wace.；Buffalo, NY: Christian Literature Publishing Co.,；1890.；<http://www.newadvent.org/fathers/3702f.htm>.}
